\documentclass[11pt]{article}
\usepackage{amsmath,amssymb,amsthm,mathtools}
\usepackage[margin=1in]{geometry}
\usepackage{booktabs}
\usepackage{hyperref}

\newtheorem{theorem}{Theorem}
\newtheorem{lemma}[theorem]{Lemma}
\newtheorem{proposition}[theorem]{Proposition}
\newtheorem{corollary}[theorem]{Corollary}
\theoremstyle{definition}
\newtheorem{remark}[theorem]{Remark}

\newcommand{\wt}[1]{\widetilde{#1}}
\renewcommand{\Re}{\operatorname{Re}}
\newcommand{\ord}{\operatorname{ord}}
\newcommand{\Sym}{\operatorname{Sym}}
\newcommand{\Sha}{\mathrm{Sha}}

\title{Carrier, Fiber, and Clock:\\
A Harmonic Transport Framework for Beyond Endoscopy and Functoriality}
\author{}
\date{}

\begin{document}
\maketitle

\begin{abstract}
Several estimates on the analytic side of the trace formula are correct \emph{bookkeeping} whose
\emph{mechanism} is left outside the frame. We isolate two. The \emph{product law}: a
derivative/repulsion statistic is identically a product of distances between configuration points, so
``repulsion'' is an algebraic identity shared by random matrices and the arithmetic. The \emph{exact
gauge}: an oscillatory integrand is a real magnitude times a deterministic phase carrier, so its
oscillation is a removable chart artifact and $|\int g|\le\int|g|$ is lossless.

Working the exact gauge out on the central analytic obstruction of Altu\u{g}'s execution of Beyond
Endoscopy for $\mathrm{GL}(2)$ --- Proposition~5.2 of Part~III, the ``central issue'' of Part~II ---
gives the paper's four contributions: (1) an \emph{oscillatory cancellation} reducing the uniformity
to a single magnitude bound with no oscillation estimate; (2) the identification of the difficulty as
a deterministic \emph{emergent clock}, the edge-comb of the fixed cutoff; (3) the consequent
\emph{removal of the productivity boundary}, re-read as gradual attrition rather than a wall; and (4)
a \emph{candidate automorphic object} in the primary representation, whose four converse-theorem
inputs --- functional equation, continuation, exhaustion, and loss-ledger recovery --- are formalized
and machine-checked.

The same two mechanisms organize three further objects: the GUE-level derivative statistics of zeros,
the archimedean transfer factor and arithmetic transparency of symmetric-power functoriality, and a
detection identity we close on a genuine transfer --- on a CM form the pole order of a self-twist
Rankin--Selberg $L$-function, read as the zero-frequency residue of a logarithmic derivative, equals a
geometric self-twist relation on Frobenius, both giving the transferred count $1$; on a non-CM form
both give $0$. The uniformity statements are theorems; the detection identity rests on classical
Gelbart--Jacquet/Rankin--Selberg theory; and the candidate object supplies the converse-theorem inputs
for every $r$. Its \emph{local} reflection algebra --- the self-reciprocal local factor, determinant
ledger, and reflection axis --- is machine-checked in Lean at the general finite duality-stable interface
the converse theorem consumes (not only the Dirichlet model). The completed functional equation is the
carrier's own geometry --- the \emph{natural} reflection, proven invariant under the dual-compatible warps,
with a tunable FE clock supplying the adjustment wherever a warp or non-self-duality would perturb the
closure, so FE closure never rests on universal warp-invariance; and entireness, continuation, and
vertical-strip growth are the outputs of a warped Abel summation in the carrier gauge --- a Riemann--Hecke
continuation of the completed carrier theta. All three follow from one carrier-cell theorem: the adapted
unit-modulus warp has exact zero mean over each complete native cell, which discharges simultaneously the
vanishing of the residual constant mode (entireness), the bounded warped primitive (continuation), and the
polynomial vertical growth. That cell-closure --- non-circular, Galois-free, so it applies to Maass forms --- is not a fiber-uniform
theorem but a per-fiber property of each carrier$+$fiber pair, effected by the fiber's own adapted warp.
For the class of finite duality-stable harmonic fibers it is \emph{proved in Lean} (each root-of-unity
weight channel cancels over its cell), the Dirichlet family and $\chi_3$ included. For the cuspidal twists a
converse theorem consumes, entireness is the vanishing of the residual mode, which equals the residue of the
convolution at $s=1$ and vanishes exactly when $\sigma^{\vee}$ is \emph{not a constituent} of the candidate
$\Sym^r\pi$; this is unconditional for $r\le4$ (where $\Sym^r\pi$ is a known cuspidal representation) and,
for $r\ge5$, an explicitly isolated carrier-indecomposability hypothesis --- the lower-degree constituent
exclusion a converse theorem exists to certify --- not a step the local carrier data shortcuts.
\end{abstract}

\tableofcontents

\section{Two mechanisms}\label{sec:two}

The present framework arose not from a proposed alternative to functoriality, but from a root-cause
analysis of the loss of productivity in the direct Beyond Endoscopy route. Once the induced clock is
separated, the higher-order analysis becomes compositional; the resulting transport interface is then
extracted from the structure actually consumed by that composition.

Two devices recur in the analytic theory as \emph{answers} while their \emph{explanations} stay
outside the frame.

\paragraph{The product law (self-interference).} The unitary ensembles predict the local statistics
of the zeros of an $L$-function correctly, but the ensemble contains no mechanism: ``eigenvalue repulsion'' is
postulated. Yet the derivative statistic that encodes repulsion is an identity: for a polynomial
$P=\prod_j(X-r_j)$ and a simple root $\mu$,
\begin{equation}\label{eq:prod}
  P'(\mu)=\prod_{r_j\neq\mu}(\mu-r_j),
\end{equation}
so $|P'(\mu)|$ \emph{is} the product of distances from $\mu$ to the other configuration points. The
same product governs the reopening rate of an arithmetic fiber at a zero. Hence a ``GUE-level''
prediction and its arithmetic counterpart compute one object; they differ only through the
configuration fed to \eqref{eq:prod}. Repulsion is the product law, not a force. (This is an
algebraic identity; it does not by itself prove that zeros follow GUE.)

\paragraph{The exact gauge (removable oscillation).} An oscillatory integral $\int g$ is normally
attacked by estimating cancellation. But if $g=|g|\,e^{i\phi}$ with $\phi$ a \emph{deterministic}
phase carrying no arithmetic content, then the oscillation is a choice of chart and
$|\int g|\le\int|g|$ is \emph{lossless}. We fix the term once: \emph{lossless} here means \emph{with
respect to arithmetic/parameter information} --- the triangle inequality is in general strict, but the
magnitude it discards is exactly the deterministic carrier phase $e^{i\phi}$, carrying no arithmetic
cancellation; so no $(C,D)$- or $\xi$-dependent content is lost, only the chart. Choosing the gauge that
makes the integrand real turns an oscillation estimate into a magnitude bound; over a root-of-unity
carrier (the $\zeta_6$ step $2\pi/6$, which closes) the gauged object is a real harmonic.

\paragraph{Part I: Altu\u{g}'s uniformity problem.} Altu\u{g}'s execution of Beyond Endoscopy for
$\mathrm{GL}(2)$ \cite{AltI,AltII,AltIII} reduces the standard-representation trace to an archimedean
analysis whose primary difficulty is a \emph{uniformity problem}: the asymptotic expansions of the
Fourier transforms of the orbital integrals must hold with implied constants independent of every
parameter --- the $C,D$-independence that is the content of the paper's long technical appendix
(Altu\u{g}~II, Appendix~A) and that secures the power saving
$\sum_{n<X}\mathrm{tr}\,T_k(n)=O(X^{31/32+\varepsilon})$ isolating the standard representation
(Altu\u{g}~III, Thm~1.1). A second obstruction sits beyond it: Poisson summation on the $n$-sum
``works well for the standard representation and the symmetric square, however stops being productive
for higher symmetric powers'' (Altu\u{g}~III, fn.~5, after Sarnak \cite{Sar}). Part~I
(\S\S\ref{sec:gauge}--\ref{sec:sarnak}) resolves the uniformity problem by the exact gauge --- one
magnitude bound, no oscillation estimate (\S\ref{sec:gauge}) --- having first identified why it looked
hard: the difficulty is the deterministic edge-comb of the fixed cutoff, an \emph{emergent clock}
(\S\ref{sec:comb}); and it re-reads the productivity boundary as gradual attrition, not a wall
(\S\ref{sec:sarnak}).

\paragraph{Part II: functoriality and a candidate object.} Part~II
(\S\S\ref{sec:gue}--\ref{sec:converse}) turns to symmetric-power functoriality --- the statement that
$\Sym^r\pi$ is automorphic on $\mathrm{GL}(r+1)$ --- and provides a concrete automorphic object to
realize it. The two mechanisms organize the transfer (\S\ref{sec:gue}); a carrier/fiber rescaling
sends a pole to a closed cell (\S\ref{sec:carrier}); the transfer's arithmetic is transparent for the
whole even family (\S\ref{sec:fl}); and a detection identity closes on a genuine example
(\S\ref{sec:detect}). The \emph{candidate automorphic object} (\S\ref{sec:construct}) --- carrier,
fiber, and warp --- carries the local reflection and carrier-interface inputs as machine-checked theorems;
the completed global twisted assembly (functional equation, warped-Abel continuation and vertical bound,
Riemann--Hecke entireness) is isolated in \S\S\ref{sec:machinery}--\ref{sec:converse}, and
\S\ref{sec:machinery}--\S\ref{sec:converse} assemble them into a functoriality theorem whose
converse-theorem hypothesis the carrier supplies for \emph{every} $r$: classically at
$\Sym^2,\Sym^3,\Sym^4$, and by the carrier's own construction beyond, where Langlands--Shahidi stops.

\section{The exact gauge on the archimedean uniformity}\label{sec:gauge}

Altu\u{g} \cite{AltI,AltII,AltIII} executes Beyond Endoscopy for $\mathrm{GL}(2)$ and the standard
representation. Its analytic core is a uniformity estimate: after the arithmetic weights $L(1,\chi_D)$
are expanded by an approximate functional equation and the trace variable is Poisson-summed, one must
bound, uniformly in $(\xi,\nu,l,f,X)$, the Fourier transform $J_{l,f}(\xi,\nu,X)$ of a product of a
fixed cutoff $G$ and a smoothed orbital transform (Prop.~5.2, Part~III); Altu\u{g} calls the
$C,D$-independence of the implied constant ``the central issue'' \cite[App.~A]{AltII}. The transform
whose asymptotics drive the appendix is (Part~II, Def.~A.6)
\begin{equation}\label{eq:A6}
  A^{\tau,\pm}_{h_a,m}(\Phi)(x)=\frac{1}{2\pi i}\int_{(\tau)}\wt\Phi(u)\,c^{\pm}_m(u/2)\,
    \Gamma\!\Big(m+1+\tfrac{a+u}{2}\Big)\,x^{-u/2}\,du,\qquad \tau>0,
\end{equation}
for a singular profile $h_a(x)=|1-x^2|^{a/2}h_1(x)$ ($h_1$ smooth, $a>-1$) and Schwartz $\Phi$; the
parameters enter only through $x=\mp C^2D$ and a unimodular prefactor.

\begin{lemma}[the gauge bound]\label{lem:mag}
For every $\tau>0$, $m\in\mathbb N$, and Schwartz $\Phi$,
\[
  \big|A^{\tau,\pm}_{h_a,m}(\Phi)(x)\big|\ \le\ B_{\tau,m,a}(\Phi)\,x^{-\tau/2},
  \qquad
  B_{\tau,m,a}(\Phi):=\frac{1}{2\pi}\!\int_{-\infty}^{\infty}\!
     \Big|\wt\Phi(\tau+iv)\,c^\pm_m\!\big(\tfrac{\tau+iv}{2}\big)\,
     \Gamma\!\big(m{+}1{+}\tfrac{a+\tau}{2}{+}\tfrac{iv}{2}\big)\Big|\,dv,
\]
with $B_{\tau,m,a}(\Phi)<\infty$ \emph{independent of $x$} (hence of $C,D$).
\end{lemma}

\begin{proof}
On $u=\tau+iv$ one has $|x^{-u/2}|=x^{-\tau/2}$, so the triangle inequality gives the bound and $B$
contains no $x$. For finiteness: by Stirling $|\Gamma(\sigma+iv/2)|\asymp e^{-\pi|v|/4}$; the phase
$i^{\,1+m+(a+u)/2}$ has magnitude $e^{-\pi v/4}$; and $|\wt\Phi(\tau+iv)|\ll e^{-\pi|v|/2}$. The
product decays like $e^{-\pi|v|}$ as $v\to+\infty$ and $e^{-\pi|v|/2}$ as $v\to-\infty$ (there the
$\Gamma$- and phase-magnitudes cancel and $\wt\Phi$ carries the decay); the $\Gamma$-poles lie at
$u<0$ and are never met for $\tau>0$.
\end{proof}

The step $|\int g|\le\int|g|$ is lossless here: in the coordinate \eqref{eq:A6} the integrand is a
real magnitude times a deterministic phase, so no arithmetic cancellation is discarded, and the
$C,D$-dependence factors out as a pure power. There is no oscillation estimate. The bound is
numerically certified on a representative Gaussian $\Phi$: $B_{1,0,1}(\Phi)=0.686$, and the certificate
margin $\sup_x |A(x)|\,x^{\tau/2}/B=0.914\le1$ holds across $x\in[0.05,2500]$ and
$(\tau,m)\in\{0.5,1,2,4\}\times\{0,2\}$. (The certificate fixes one $\Phi$; the bound of
Lemma~\ref{lem:mag} is general, holding for every Schwartz $\Phi$ through the finite line integral
$B_{\tau,m,a}(\Phi)$.)

\begin{theorem}[the sharp expansion, Part~II Thm A.14]\label{thm:A14}
With $h_a,\Phi$ as above, for every $M$,
\[
  I(C,D):=\int_{-1}^{1}h_a(x)\,\Phi\!\big(C/\sqrt{1-x^2}\big)\,e(xD)\,dx
  =\sum_{m=0}^{M}\frac{e(\pm D)\,A^{\pm}_m(\mp C^2D)}{(\mp 2\pi iD)^{\,m+1+a/2}}+R_M(C,D),
\]
with $A^\pm_m$ as in \eqref{eq:A6} and $|R_M|\le\mathcal B_{M,\tau_1}(C^2D)^{-\tau_1/2}
D^{-(M+1+a/2)}$, $\mathcal B_{M,\tau_1}$ independent of $C,D$.
\end{theorem}

\begin{proof}
\emph{Localization.} Fix a smooth partition $1=\chi_++\chi_0+\chi_-$ with $\chi_\pm$ supported in
$|x\mp1|<\delta$ and $\chi_0$ in $[-1+\delta,1-\delta]$. On $\chi_0$ the phase $x\mapsto x$ has
nonvanishing derivative, so the non-stationary phase lemma \cite[Thm.~7.7.1]{Hormander} ($N$-fold
integration by parts) gives an $O_N(D^{-N})$ contribution with implied constant controlled by
$\|\partial_x^N(\chi_0 h_a\,\Phi(C/\sqrt{1-x^2}))\|_{L^1}$. This is bounded uniformly in $C$: on
$[-1+\delta,1-\delta]$ the factor $\Phi(C/\sqrt{1-x^2})$ and its $x$-derivatives are bounded uniformly
in $C$, since for Schwartz $\Phi$ and $c\ge\delta'>0$ one has $\sup_{C>0}|C^k\Phi^{(k)}(Cc)|<\infty$.

\emph{Endpoint reduction.} Near $x=1$ substitute $x=1-t$, $t\in(0,\delta)$: then
$h_a(1-t)=t^{a/2}\varphi(t)$ with $\varphi(t)=(2-t)^{a/2}h_1(1-t)\in C^\infty[0,\delta)$;
$\sqrt{1-x^2}=\sqrt t\,\sqrt{2-t}$; and $e(xD)=e(D)e(-tD)$, so
$I_+=e(D)\int_0^\delta \chi_+(1-t)\,t^{a/2}\varphi(t)\,\Phi\big(C/(\sqrt t\sqrt{2-t})\big)e(-tD)\,dt$.

\emph{Mellin representation and closed-form $t$-integral.} For $\sigma>0$ in the strip of holomorphy
of $\wt\Phi$, insert $\Phi(w)=\frac1{2\pi i}\int_{(\sigma)}\wt\Phi(u)w^{-u}\,du$ with
$w=C/(\sqrt t\sqrt{2-t})$. The joint integrand is absolutely integrable on
$\{\Re u=\sigma\}\times(0,\delta)$ --- the $t$-integral converges absolutely and $\wt\Phi(\sigma+iv)$
decays exponentially in $v$ --- so Fubini applies. Writing
$\psi_u(t)=\varphi(t)(2-t)^{u/2}\chi_+(1-t)\in C^\infty$ and Taylor-expanding
$\psi_u(t)=\sum_{m=0}^{M}b_m(u)t^m+R_M(u,t)$, the $m$-th $t$-integral is, by Watson's lemma / the
Mellin evaluation of an oscillatory monomial \cite[Ch.~3]{Wong} (the smooth cutoff at $t=\delta$ costs
only $O(D^{-\infty})$),
\[
  b_m(u)\!\int_0^\infty t^{m+\frac{a+u}{2}}e(-tD)\,dt
  = b_m(u)\,\Gamma\!\big(m{+}1{+}\tfrac{a+u}{2}\big)(2\pi iD)^{-(m+1+\frac{a+u}{2})}+O(D^{-\infty}).
\]
The oscillation is evaluated in closed form, with no cancellation estimate. Since
$C^{-u}D^{-u/2}=(C^2D)^{-u/2}$, collecting the $u$-contour integral gives the $m$-th term
$e(D)(2\pi iD)^{-(m+1+a/2)}A^+_m(C^2D)$ with $A^+_m$ exactly \eqref{eq:A6}: the coefficients $b_m(u)$,
carrying the $(2-t)^{u/2}$ Taylor data, reproduce Altu\u{g}'s $c^\pm_m$. The identification is forced, not
matched by hand: both $b_m(u)$ and $c^\pm_m(u/2)$ are the coefficient of order $D^{-(m+1+a/2)}$ in the
asymptotic expansion of the \emph{same} endpoint integral $I_\pm(C,D)$ in the \emph{same} asymptotic scale
$\{D^{-(m+1+a/2)}\}_{m\ge0}$ --- Altu\u{g}'s by his Watson-lemma reduction, ours by the Mellin evaluation
above --- and an asymptotic expansion in a fixed scale is unique \cite[Ch.~1]{Wong}: the coefficient
functionals $\lim_{D\to\infty}D^{\,m+1+a/2}\big(I_\pm-\sum_{j<m}(\cdots)\big)$ are determined by $I_\pm$ and
the scale alone, so $b_m(u)=c^\pm_m(u/2)$ termwise.

\emph{Remainder.} The Taylor remainder satisfies $|R_M(u,t)|\le C_M(u)\,t^{M+1}$ with
$C_M(u)=\sup_{[0,\delta]}|\partial_t^{M+1}\psi_u|/(M+1)!$ of at most polynomial growth in $\Im u$
($\psi_u$ smooth). Its contribution is again a transform of the form \eqref{eq:A6} with amplitude
$R_M$; taking its $u$-contour at $\Re u=\tau_1$ (any $\tau_1>0$), Lemma~\ref{lem:mag} bounds it by
$\mathcal B_{M,\tau_1}(C^2D)^{-\tau_1/2}D^{-(M+1+a/2)}$ with $\mathcal B_{M,\tau_1}<\infty$
independent of $C,D$. The endpoint $x=-1$ gives the $e(-D)$, $A^-_m$ terms identically.
\end{proof}

\begin{remark}[on the exponents]
The contour height $\tau_1>0$ is free, and the remainder decays like $(C^2D)^{-\tau_1/2}$ because the
factor $(C^2D)^{-u/2}$ is evaluated at $\Re u=\tau_1$; Altu\u{g}'s stated $(C^2D)^{-\tau_1}$ is the same
statement in his normalization of the Mellin variable ($u$ versus $u/2$). What the application uses is
only that, for arbitrary $A>0$, the remainder is $O_A\big((C^2D)^{-A}D^{-(M+1+a/2)}\big)$ uniformly in
$C,D$ --- immediate from the freedom in $\tau_1$ together with Lemma~\ref{lem:mag}.
\end{remark}

\begin{proposition}[Part~III Prop.~5.2]\label{prop:52}
$|J_{l,f}(\xi,\nu,X)|\ll X^2\nu^{-M}(lf^2)^{-N}\big((lf^2/\sqrt X)^{N-M+3}+\xi^M\big)$ with implied
constant independent of $(\xi,\nu,l,f,X)$.
\end{proposition}

\begin{proof}
Write $J_{l,f}(\xi,\nu,X)=\int G(y/X)\,y\,\Psi(y)\,e(-y\nu/2lf^2)\,dy$ with $\Psi(y)=I_{l,f}(\xi,y)$.
Since $G$ is compactly supported, integrating by parts $M$ times against the phase gives the exact
identity $J=(lf^2/(-\pi i\nu))^M\int\partial_y^M(G(y/X)\,y\,\Psi(y))\,e(-y\nu/2lf^2)\,dy$, hence
\[
  |J_{l,f}(\xi,\nu,X)|\le\Big(\tfrac{lf^2}{\pi\nu}\Big)^{M}
    \big\|\partial_y^M\!\big(G(y/X)\,y\,\Psi(y)\big)\big\|_{L^1(y)} .
\]
By the Leibniz rule
$\partial_y^M(G(y/X)\,y\,\Psi)=\sum_{i+j+k=M}\binom{M}{i,j,k}X^{-i}G^{(i)}(y/X)\,(y)^{(j)}\,
\partial_y^k\Psi$, where $(y)^{(j)}=y,1,0$ for $j=0,1,\ge2$. The cutoff $G^{(i)}(y/X)$ is supported on
$y\in[\tfrac X4,\tfrac{5X}4]$ (length $\asymp X$) and contributes $\|G^{(i)}\|_\infty$, and the factor
$(y)^{(0)}=y\asymp X$ supplies a second power of $X$, giving the $X^2$. Each $y$-derivative of
$\Psi(y)=\int\theta_\infty(x)[F+H](\arg)\,e(-x\xi\sqrt{4y}/4lf^2)\,dx$ falls on either the phase ---
producing a factor $O(\xi/lf^2\sqrt y)$, the source of the $\xi^M$ branch --- or on the profile through
$\arg=lf^2/(\sqrt{4y}\sqrt{1-x^2})$, producing a factor $O(lf^2/y)$, the source of the $(lf^2/\sqrt
X)$ branch. Each resulting $x$-integral is a transform of the form \eqref{eq:A6} (with $C\asymp\xi\sqrt
y/lf^2$ and $D$ the local scale), bounded by Lemma~\ref{lem:mag}, or at the sharp order by the
expansion of Theorem~\ref{thm:A14}. Collecting the $X$-, $lf^2$-, and $\xi$-powers reproduces the two
branches $(lf^2/\sqrt X)^{N-M+3}$ and $\xi^M$; the implied constant is a finite sum of the integrals
$B_{\tau,m,a}(\Phi)$ (Lemma~\ref{lem:mag}) and the norms $\|G^{(i)}\|_\infty$, none depending on
$(\xi,\nu,l,f,X)$.
\end{proof}

Every uniform constant in \S\ref{sec:gauge} is one application of Lemma~\ref{lem:mag}; the content of
Altu\u{g}'s Appendix~A, from the exact-gauge standpoint, is that single magnitude estimate. Where
Altu\u{g}'s route controls the transform through the asymptotic expansion together with an oscillation
bound on the remainder, the exact gauge replaces that oscillation bound by the triangle inequality in
the coordinate that makes the integrand real: the whole $(\xi,\nu,l,f,X)$-uniformity then rests on the
$x$-independence of the single line integral $B_{\tau,m,a}(\Phi)$, with no cancellation estimate at any
step.

\section{The oscillation is the cutoff's comb}\label{sec:comb}

Why did Prop.~\ref{prop:52} look hard? The oscillation the estimate must survive is deterministic.
The numerical scan is made in the clock coordinate
\[
  \nu_{\mathrm{clock}}:=2\nu_J,\qquad
  e(-y\nu_J/2lf^2)=e(-y\nu_{\mathrm{clock}}/4lf^2),
\]
where $\nu_J$ is the transform frequency in Prop.~\ref{prop:52}. The measured clock is extracted from the
full complex profile $V(\nu_{\mathrm{clock}})$ by fitting the log-power modulation
$\log |V(\nu_{\mathrm{clock}})|^2$ to a smooth envelope plus the first two harmonics sharing one
fundamental, not by averaging dip gaps. Over the cutoff support $y\in[\tfrac X4,\tfrac{5X}4]$ this gives
\[
  \Delta\nu_{\mathrm{clock}}
   =0.524\,(lf^2)^{1.03}(X/8)^{-1.04}(1+\xi)^{0.00},\qquad R^2=1.000,
\]
equivalently half that spacing in the $\nu_J$ coordinate. In the one-parameter window-span normalization
$\kappa_{\mathrm{eff}}:=4lf^2/(X\Delta\nu_{\mathrm{clock}})$ the cells give
$\kappa_{\mathrm{eff}}=0.943\pm0.025$ (range $0.895$--$0.971$), while the harmonic-truncation audit has
median period shift $2.4\times10^{-4}$, so the residual spread is not a dip-localization artifact. With the
same Chebyshev multiplier $U_r(x)$ used for the symmetric-power rank, the event-bearing $\xi\ne0$ cells keep
the same law through the tested range $0\le r\le13$ after separating the continuous clock from deep-event
visibility: every rank has $R^2=1.000$, $lf^2$-exponent $1.03$--$1.04$, and $X$-exponent
$-1.04$--$-1.03$ on the cells passing the clock-quality gate. Thus the older scalar
$\kappa\approx0.94$ is the mean of a slightly drifting clock, and the observed high-rank loss is productivity
decay: deep-event cells thin from $10/12$ at $r=0$ to $3/9$ and $3/8$ at $r=11,12$, while the surviving
clock cells still obey the same law. This is
the interference of the two edges of the fixed cutoff, switched on once the $\xi$-phase drives the $x$-integral
into edge dominance. At $\xi=0$ the $x$-integrand is unimodal (prominence ratio $P_2/P_1=0$) and the deep-dip
comb is absent. The apparent difficulty is a chart feature of
the fixed cutoff---exactly what the exact gauge of \S\ref{sec:gauge} removes.

\section{The productivity boundary}\label{sec:sarnak}

The symmetric-power kernel is modulated by $U_r(\cos\theta)=\sum_{j=0}^r e^{i(r-2j)\theta}$, carrying
a zero-frequency component iff $r$ is even. At $\xi=0$ the $x$-integral reads this component (the
detecting residue); at $\xi\neq0$ it is the removable comb of \S\ref{sec:comb}, bounded by
Prop.~\ref{prop:52} (with $|U_r|\le r+1$) and hence $o(X)$. The two regions are disjoint in $\xi$:
the detecting residue sits at $\xi=0$, where the comb is absent; the comb sits at $\xi\neq0$, where it
is removable. So the residue isolates for every $r$. On the dual integral (level $p^k=25$) the
zero-frequency value against the $\xi\neq0$ floor is
\[
  r=1:\ 0\ (\text{exact}),\quad r=2:\ 55.9,\quad r=3:\ 1.4\times10^{-16},\quad r=4:\ 11.4,
\]
nonzero iff $r$ even, clear of the floor by an order of magnitude. The degradation with $r$ is the
growth of the character tail mass ($0.042,0.092,0.135,0.156,0.256$ for $r=0,\dots,4$, fit
$\text{ceiling}\times\sqrt{\#\text{harmonics}}$ at $R^2=0.977$)---gradual attrition, not a wall.
Sarnak's productivity boundary \cite[fn.~5]{AltIII} records the growth of a removable floor, not a
loss of the residue.

\section{Toward functoriality: the product law and the transfer factor}\label{sec:gue}

The three preceding sections close the analytic core --- the uniformity (\S\ref{sec:gauge}), the
emergent clock (\S\ref{sec:comb}), and the productivity boundary. The remainder of the paper develops
the two mechanisms toward symmetric-power functoriality, beginning with the product law, which supplies
the transfer factor.

The derivative statistic that random-matrix theory reads at an eigenvalue is the reopening rate an
arithmetic fiber reads at a zero, by \eqref{eq:prod}: both are $\prod_{r_j\neq\mu}|\mu-r_j|$. This is
an identity, not a coincidence, and it is the mechanism the ensemble leaves implicit. The consequence
here is structural: the same distance product is a \emph{Weyl denominator}, and Weyl denominators are
the archimedean transfer factors of \S\ref{sec:fl}. Self-interference is one object appearing as a
zero statistic, a reopening rate, and a transfer factor.

\begin{remark}[a distinct post-closure mechanism, deferred]\label{rem:reverb}
The product law fixes the derivative statistic \eqref{eq:prod} but not the \emph{configuration} that
feeds it, and that configuration is governed by a separate dynamics we only flag here. At a focal
cancellation the ordinary phasor residue vanishes, yet the fiber carries a nonzero eigen response that
persists over roughly one native harmonic cell; when closures fall within a cell their response
intervals overlap and deform the local spacing deterministically. The algebraic backbone --- a
\emph{cluster product law} generalizing \eqref{eq:prod} to any finite cluster --- is proven
(\texttt{ReverbResidue.cluster\_product\_law}, with pair suppression and the reopening residue), and
the dynamics is measured to hold uniformly across pairs, triples, quads, and clusters carrying
non-uniformly spaced neighbours: it is not a pairwise artifact, the case a triple would expose. This is
a deterministic candidate for the configuration law that populates \eqref{eq:prod} --- like the
unitary ensembles, but with a reason. We make no limiting-ensemble claim here (finite-$N$
characteristic-polynomial geometry is naturally unitary-circular, while the unfolded zero statistics
live in the limiting language); formulating the response kernel and deriving the unfolded $n$-level
correlations is the subject of a separate paper.
\end{remark}

\section{The carrier/fiber rescaling: a pole becomes a closed cell}\label{sec:carrier}

Place the integers on a carrier scaled by $\pi/3$, i.e. attach the cell phase $e^{in\pi/3}$ (the
$\mu_6$/$\zeta_6$ step). We distinguish two operations on the carrier throughout. A \emph{rescaling} is a
\emph{fixed constant} --- here $\pi/3$, or another $\mu_m$ step --- that lifts a unit-$1$ carrier into one
with harmonic cells; it \emph{establishes} the native cell structure (and is the reason we never work at
unit scale). A \emph{warp} is a \emph{function}, a continuously-evolving deterministic output
$\omega(\cdot)$ that rides the already-rescaled carrier and adapts to the fiber --- the $\omega_W(n)$ of the
warped-Abel summation (\S\ref{sec:construct}) and the tunable FE clock (\S\ref{sec:general}). Rescaling sets
the cell; warping moves within it. The pole of $\zeta$ at $s=1$ lives in the divergence of $\sum 1/n$; on
the (rescaled) carrier it closes,
\[
  \sum_{n\ge1}\frac{e^{in\pi/3}}{n}=-\log\!\big(1-e^{i\pi/3}\big)=\frac{i\pi}{3},
\]
a finite $\mu_6$ value: the divergence has become a closed-cell reading. For a Dirichlet character
$\chi$ of conductor $q$, the appropriate fiber scale (trivial $\to\pi/6$, $\chi_3\to\pi/3$,
$\chi_4\to\pi/2$) makes carrier and fiber compatible harmonics. The character $\chi_3$ is the
\emph{aligned archetype}: its conductor $3$ is the Eisenstein conductor --- the sixth roots of unity of
$\mathbb Z[\zeta_3]=\mathbb Z[\zeta_6]$ --- so the $\pi/3=\mu_6$ rescaling coincides with the $L$-function's
own conductor symmetry, the carrier cell \emph{is} the fiber's native cell, and closure needs the rescaling
alone with no warp: hence \emph{no residue}, the residual-mode vanishing of Lemma~\ref{lem:residual} in its
cleanest, rescaling-only instance. The principal character, carrying the
pole, is handled by the eta regulator,
\[
  \sum_{n\ge1}\frac{(-1)^{n-1}e^{in\pi/3}}{n}=\log\!\big(1+e^{i\pi/3}\big)=\tfrac12\ln3+\tfrac{i\pi}{6},
\]
landing on the trivial character's own cell ($\pi/6$, $\mu_{12}$), pole-free, with the partial-sum
phase converging (no drift). In this frame the $S(t)$ oscillation of the readout is removed---the
value is a fixed harmonic, not a wandering phase---so what was a residue to be identified is a
zero-frequency cell value, read as the DC-residue $(s-c)F'/F\to r$ of a logarithmic derivative. This
is the reading used in \S\ref{sec:detect}.

\section{Symmetric-power transfer: arithmetic transparency}\label{sec:fl}

For $\Sym^r$ transfer $\mathrm{GL}(2)\to\mathrm{GL}(r+1)$ ($r=2m$ even), the transfer sends an
elliptic class of Satake angle $\theta$ to the class of angles $\{2k\theta:|k|\le m\}\cup\{0\}$
---i.e. $\Sym^r$ of the Frobenius block, with the fixed eigenvalue $1=\alpha\beta$ the det-one of
the block. Two features make the archimedean transfer factor explicit and the arithmetic transparent.

\emph{Arithmetic transparency.} The transferred characteristic polynomial is reducible ($x-1$
divides it), and each quadratic factor $x^2-2\cos(2k\theta)x+1$ has discriminant $-4\sin^2(2k\theta)$.
The clock identity $\sin(2k\theta)=\sin\theta\,U_{2k-1}(\cos\theta)$ (with $\sin^2\theta=|D|/4n$,
$D=t^2-4n$) gives
\[
  \frac{\operatorname{disc}_k}{D}=\frac{U_{2k-1}(\cos\theta)^2}{4n}\in\mathbb Q^{\times2},
\]
a rational square---so every factor lies in the \emph{same} field $\mathbb Q(\sqrt D)$ as the
original class, the arithmetic $L$-value is identical on both sides, and it cancels in the transfer
factor. (Because symmetric powers transfer to \emph{reducible} classes, this avoids the order-zeta
conjecture of Deng--Espinosa, which concerns irreducible cubic orders.)

\emph{Transfer factor as a distance product.} What remains is the ratio of Weyl denominators
---distance products by \eqref{eq:prod}. For $r=2$,
\[
  \frac{\Delta_3}{\Delta_2}
  =\frac{|1-e^{2i\theta}|^2\,|e^{2i\theta}-e^{-2i\theta}|}{|e^{i\theta}-e^{-i\theta}|}
  =4|\sin\theta||\sin2\theta|=8\sin^2\theta|\cos\theta|=\frac{|D|\,|t|}{n^{3/2}},
\]
a monomial in the class data; its winding phase, gauged by the $\pi/3$ ($\mu_6$) step, is a real
harmonic (\S\ref{sec:carrier}). Thus $\Delta(\theta)=U_r(\cos\theta)\cdot(\Delta_{r+1}/\Delta_2)
\cdot(\text{radial }\sqrt p\text{-scaling})$, the Weyl character and distance product being product-law
objects and the phase an exact-gauge object.

\section{Detection: the self-twist identity closes}\label{sec:detect}

We isolate one computable instance of the detection step and close it on a genuine transfer. For a
holomorphic newform $f$ with $a_p$ its Frobenius trace and $\lambda(p)=a_p/p^{(k-1)/2}$, the
symmetric square carries the middle eigenvalue $\alpha\beta=1$; the plain census
$\sum_p(\lambda(p)^2-1)(\log p)/p$ has mean zero and is blind. The correct object is the
\emph{self-twist} Rankin--Selberg $L$-function $L(s,f\times f\otimes\chi_K)$ for the quadratic
character $\chi_K$ of an imaginary quadratic field $K$: it has a pole at $s=1$ iff $f$ has complex
multiplication by $K$ (dihedral), the classical criterion of Gelbart--Jacquet \cite{GJ} and
Rankin--Selberg.
Its pole order is the transferred count, read as the zero-frequency residue of the logarithmic
derivative:
\[
  \text{(LEFT, analytic)}\qquad
  A_{\mathrm{tw}}(X)=\sum_{p\le X}\lambda(p)^2\chi_K(p)\frac{\log p}{p},\qquad
  \frac{d\,A_{\mathrm{tw}}}{d\log X}\ \longrightarrow\ \ord_{s=1}L(s,f\times f\otimes\chi_K).
\]
The geometric side is the self-twist relation on Frobenius, $f\otimes\chi_K\cong f$, equivalently
$a_p=0$ at every inert prime:
\[
  \text{(RIGHT, geometric)}\qquad
  \#\{p\le X\ \text{inert}:\ a_p=0\}\big/\#\{p\le X\ \text{inert}\}\ \longrightarrow\ 1_{\{f\ \text{has CM by }K\}}.
\]

We evaluate both on two weight-two forms with $K=\mathbb Q(\sqrt{-3})$ (the Eisenstein field, whose
$\chi_K(p)=+1,-1$ for $p\equiv1,2\pmod3$): the CM curve $y^2=x^3+1$ and the non-CM curve
$y^2=x^3-16x+16$ (conductor $37$), with $a_p$ from point counting.

\begin{center}
\begin{tabular}{lccc}
\toprule
 & LEFT: slope $dA_{\mathrm{tw}}/d\log X$ & RIGHT: inert $a_p=0$ fraction & count \\
\midrule
CM $\sqrt{-3}$ ($x^3+1$) & $0.997$ & $725/726=0.9986$ & $1$ \\
non-CM ($x^3-16x+16$) & $0.037$ & $6/726=0.0083$ & $0$ \\
\bottomrule
\end{tabular}
\end{center}

For the CM form the analytic slope is $0.97$--$0.99$ across every interval $[1000,4000]$,
$[4000,10000]$, $[10000,20000]$ (pole order $1$), and the geometric self-twist holds at every good
inert prime (the single exception $725/726$ is $p=2$, the bad-reduction prime of conductor
$36=2^2\cdot3^2$). For the non-CM form both sides are $0$. Thus
\[
  \boxed{\ \ord_{s=1}L(s,f\times f\otimes\chi_K)\ =\ 1_{\{f\otimes\chi_K\cong f\}}\ =\ \text{transferred count}\ }
\]
closes---$1$ on a genuine transfer, $0$ on a genuine non-transfer---read through the logarithmic-
derivative DC-residue on the S(t)-free carrier of \S\ref{sec:carrier}, with the emergent-clock
lesson explaining why the plain census cancels (inert primes carry $\lambda^2=0$, so the self-twist
census has no cancelling floor).

\begin{remark}[what this is]
The mathematics detected here---that the symmetric-square/self-twist pole marks CM forms---is
classical (Gelbart--Jacquet, Rankin--Selberg). What the section contributes is the closure of the
detection \emph{identity} through the two mechanisms: the analytic count as a product-law/log-
derivative residue, the geometric count as the self-twist relation, matched on a genuine transfer.
\end{remark}

\subsection{A non-abelian instance: which Frobenius data transfers}\label{sec:a4}

The self-twist detection of \S\ref{sec:detect} runs on a dihedral (abelian, $C_2$) form, where the
Galois image is small and the arithmetic reduces to a single quadratic character. For a genuinely
non-abelian form the Frobenius datum splits into two independent bits---one a character cannot
see---and the transfer treats them differently. We illustrate this on the exotic tetrahedral
(``$A_4$'') weight-one newform $h$ of level $124$, whose $L$-coefficients are a class function of
$\mathrm{Frob}_p$ in the binary tetrahedral group $2T=\mathrm{SL}(2,3)$, given by the closed formula
$a_p=\zeta_{12}^{k(p)}$ with
\[
  k(p)=\tfrac12\big(\chi_k(p)+4\varepsilon(p)\big)+\big(0\ \text{if }(3/p)=+1,\ \text{else }6\big)
  \pmod{12},
\]
$a_p=0$ when $\varepsilon(p)=0$. Here the exponent carries two independent Galois bits:
\[
  \text{\emph{value bit}}\ \varepsilon(p)\in\{0,\pm1\}\ \text{(the cubic residue, the abelian
  $A_4^{\mathrm{ab}}=C_3$ datum)},\qquad
  \text{\emph{sign bit}}\ (3/p)\ \text{(the $\mathrm{SL}(2,3)\!\to\!A_4$ central $\pm1$,}
\]
the double-cover datum read from the order of $\mathrm{Frob}_p$; the sign bit resolves $k$ versus
$k+6$ and is invisible to any abelian character.

The symmetric square sends $a_p\mapsto a_p^2-\chi(p)$, and $a_p^2=\zeta_{12}^{2k}$ is unchanged by
$k\mapsto k+6$. Hence:

\begin{center}
\begin{tabular}{lc}
\toprule
 & Sym$^2$ transfer to GL(3) \\
\midrule
sign bit $(3/p)$ (double cover, $\mathrm{SL}(2,3)$) & \emph{forgotten} (trace invariant under $k\!\leftrightarrow\!k{+}6$) \\
value bit $\varepsilon$ (cubic residue, projective $A_4$) & \emph{retained}: $\varepsilon=\pm1\Rightarrow\pm\sqrt3\,i$, $\varepsilon=0\Rightarrow\pm1$ \\
\bottomrule
\end{tabular}
\end{center}

So the transfer factors through the projective image $A_4$: the value bit passes to the GL(3) Satake
datum, the sign bit is the GL(2)-specific double-cover fiber the transfer collapses. The order trick
---sign from the element order, value from the residue symbol---thus reads \emph{exactly which part
of the Frobenius datum transfers}, and the split is non-trivial precisely because $A_4$ is
non-abelian (for the dihedral case of \S\ref{sec:detect} the two bits coincide). This identifies the
transferred datum; the automorphy of the resulting GL(3) object (the three-dimensional standard
representation of $A_4$) is the classical Artin/Langlands--Tunnell case.

\subsection{The detected count is a Mordell--Weil rank difference}\label{sec:rankdiff}

The transferred object of \S\ref{sec:a4} is, up to the nebentype twist, the three-dimensional
adjoint $\mathrm{Ad}_g=\mathrm{End}^0(V_g)\cong\Sym^2 V_g\otimes(\det V_g)^{-1}$ of the $A_4$ form
$g$. At this exotic locus the zero-frequency count read on the analytic side coincides with a
genuine arithmetic rank, and the coincidence is an exact identity of group theory.

The quartic field $M$ cut out by $g$ carries the permutation action of $A_4$ on four points, whose
character is $\chi_{\mathrm{perm}}=(4,0,1,1)$ on the classes $(e,\ (2,2),\ (3)_a,\ (3)_b)$. Against the
character table, $\langle\chi_{\mathrm{perm}},\mathbf 1\rangle=1$ and
$\langle\chi_{\mathrm{perm}},\chi_{\mathrm{std}}\rangle=1$, so
\[
  \chi_{\mathrm{perm}}=\mathbf 1\ \oplus\ \mathrm{Ad}_g,\qquad \mathrm{Ad}_g=\chi_{\mathrm{std}}\
  \text{(the $3$-dimensional irreducible)} .
\]
Functoriality of Mordell--Weil under this decomposition gives
$E(M)\otimes\mathbb Q = (E(\mathbb Q)\otimes\mathbb Q)\oplus E(\,\cdot\,)_{\mathrm{Ad}_g}$ and hence the
\emph{exact} Artin/permutation identity
\begin{equation}\label{eq:rankdiff}
  r(E,\mathrm{Ad}_g)\ =\ \operatorname{rank}E(M)-\operatorname{rank}E(\mathbb Q),
\end{equation}
proven as group theory (no analytic or $p$-adic input). The left side is the order of vanishing at
$s=1$ of $L(E,\mathrm{Ad}_g,s)$---the zero-frequency residue of the logarithmic derivative, the same
statistic \S\ref{sec:detect} read as the transferred count. The right side is a difference of
classical Mordell--Weil ranks, the source living ``up the tower'' at $M$ and traced down to
$\mathbb Q$. Thus the count is a rank difference: one statistic, read once on the automorphic side
and once on the arithmetic side.

On the two Darmon--Lauder--Rotger exotic $A_4$ curves this is pinned on both sides:
\begin{center}
\begin{tabular}{lccc}
\toprule
 & $\operatorname{rank}E(\mathbb Q)$ & $\operatorname{rank}E(M)$ & $r(E,\mathrm{Ad}_g)=\operatorname{rank}E(M)-\operatorname{rank}E(\mathbb Q)$ \\
\midrule
$26b$ ($y^2+xy+y=x^3-x^2-3x+3$) & $0$ & $2$ & $2$ \\
$52b$ ($y^2=x^3+x-10$) & $0$ & $2$ & $2$ \\
\bottomrule
\end{tabular}
\end{center}
with $M=\mathbb Q(w)$, $w^4+7w^2-2w+14=0$ the tetrahedral field of \S\ref{sec:a4}, rank-two
generators of $E(M)$ exhibited by descent. This is the ``one mechanism, two locations'' statement
made theorem-anchored: at the exotic locus both sides of \eqref{eq:rankdiff} are determined, the
identity between them is proven group theory, and the shared statistic is the DC-residue that reads
the functorial pole in \S\ref{sec:detect} and the analytic rank in the Beyond-Endoscopy census
alike.

\begin{remark}[tiering]
The rank identity \eqref{eq:rankdiff} is proven (Artin formalism). The automorphy of $\mathrm{Ad}_g$
on GL(3)---which makes $L(E,\mathrm{Ad}_g,s)$ an honest automorphic $L$-function so that its pole order
is a spectral quantity---is the classical Langlands--Tunnell case for $A_4$. The $p$-adic
\emph{regulator} side (the explicit point coordinates realizing the rank) is the Darmon--Lauder--Rotger
Elliptic Stark conjecture, verified to $20$ digits on $26b$/$52b$ but with the points found by
descent first; we use only the proven rank identity and the classical automorphy here.
\end{remark}

\section{A candidate automorphic object from the primary construct}\label{sec:construct}

The transfer of \S\S\ref{sec:fl}--\ref{sec:rankdiff} requires an automorphic representation on
$\mathrm{GL}(r+1)$ whose existence realizes the functorial lift. In the primary (three-dimensional)
representation we exhibit a concrete \emph{candidate object} for it: the \emph{carrier} --- a
double-sided helix of Frobenius similitudes --- carrying a \emph{fiber} (the phasor readout), together
with its \emph{warp} (the twist family of \S\ref{sec:fl}). The classical $L$-function is the shadow of
this candidate under the projection that sends the exponential state space $e^{y}$ down by a
M\"obius/Cayley map and then $\log$ to the ordinate, at the midpoint and without drift. What makes it
a candidate is that the three properties a converse theorem asks of an $L$-function --- functional
equation, analytic continuation, and boundedness/completeness --- are each a \emph{proven} statement
about it, and each is already formalized.

\paragraph{Functional equation: the reflection axis.} The anti-helix is the complex \emph{conjugate}
of the helix, not its mirror: the carrier's transverse block is $\operatorname{diag}(z,\bar z)$ with
$z$ the Frobenius spin. It is literally the de Branges conjugate-pair block
(\texttt{frobeniusBlock\_eq\_conjPairBlock}), with $z\bar z=|z|^2=1$ (\texttt{spin\_mul\_conj}), so
$\det=1$ and it is unitary (\texttt{frobeniusBlock\_det\_one}, \texttt{frobeniusBlock\_unitary}). The
det-one condition is not a normalization but the \emph{reality locus}:
$\det=1\iff|z|=1\iff$ the helix and anti-helix \emph{balance} $\|E(z)\|=\|E^{*}(z)\|\iff z$ real
(\texttt{frobeniusBlock\_deBranges\_reality\_bridge}). The dualizing reflection $s\mapsto1-s$ fixes
\emph{exactly} the critical line, $\Re(1-s)=\Re s\iff\Re s=\tfrac12$ (\texttt{reflection\_fixes\_iff}),
and this abscissa is not posited but \emph{forced} by the arclength area law --- the carrier admits a
scale-balanced fiber iff $\sigma=\tfrac12$ (\texttt{scaleBalanced\_iff}), the reciprocation of the
fiber amplitude $n^{-\sigma}$ against the area-law radius $\sqrt n$. This is the functional equation
realized geometrically: the conjugate leg closing on the reflection-fixed axis. It survives
\emph{every} unit-modulus warp $z\mapsto Az$ ($|A|=1$), $z\bar z\mapsto|A|^2=1$
(\texttt{warpedBlock\_det\_one\_of\_warp}), so the whole twist family carries it, structurally and
without estimate. (A complex twist is \emph{self}-dual only when it equals its own conjugate;
otherwise the equation pairs the sideband $\chi_j$ with $\chi_{-j}=\bar\chi_j$, as the conjugate leg
demands.)

\paragraph{The crossing geometry and the two fiber topologies.} A zero is a phasor \emph{focal
cancellation}: the phasors, generated continuously as the carrier grows, align and cancel at a height
fixed by the carrier scale, the order of the integers/primes, and the conductor. Because the same data
drive both conjugate legs, the two legs cancel at the \emph{same} height, the determinant-one link
holding between them. Read as its own oscillator the fiber accumulates phase along each leg; the
crossing spacing is a half-turn on each side of the origin and a full turn thereafter. This reduces to
the functional equation making the central fiber an extremum --- $Z$ real and even, $Z'(0)=0$
(\texttt{collapseWave\_even}, \texttt{hinge\_turning\_point}) --- together with the argument principle,
so the half-turn offset and full-turn step are provable; only the \emph{rigidity} of the spacing in $t$
is not settled by this argument. The radial envelope of the winding is the Archimedean spiral, whose arclength area
law $r_n\sim\sqrt n$ forces the abscissa to $\tfrac12$ (\texttt{scaleBalanced\_iff}).

Two fiber topologies occur. An \emph{open} helix --- the Dirichlet family --- has its fibers begin at
the origin and run up both conjugate legs to infinity. A \emph{clipped} helix --- a self-dual
$L$-function such as that of an elliptic curve --- has its strand bounded by the conductor: the live
phasor count scales as $\sqrt N$, the analytic conductor (measured exponent $0.5000$ across seven
decades, Appendix~\ref{app:repro}), the same $\sqrt{\,\cdot\,}$ law that fixes the abscissa; and its
two legs \emph{meet} at the central point --- the fixed point of the functional-equation involution
(UNIT/2, \texttt{affine\_reflection\_fixed\_iff}), the ordinate origin $t=0$ (equivalently
$s=\tfrac12$ analytic, $s=1$ arithmetic), not the spiral base $N=0$. There the two strands carry the
determinant-one conjugate weights $r^{s-1}\!\cdot r^{1-s}=1$ (\texttt{strand\_weights\_det\_one}), and
the root number $\varepsilon$ \emph{is} the weld sign: $\varepsilon=-1$ cancels the central term and
forces a central zero (\texttt{weld\_kills\_each\_phasor}), $\varepsilon=+1$ doubles it
(\texttt{weld\_doubles\_each\_phasor}). The same double-ended weld reads, in one factorization
$(s-c)^r G$, the root number, the rank $r$ (the DC-residue, \texttt{rank\_is\_dc\_residue}), and the
leading datum $L^{(r)}/r!$ --- matched to tabulated values on ranks $0$–$3$.

\paragraph{Continuation: the readout, continued past the strip.} The fiber's value is the phasor
representation $L(\chi,s)=\sum_n\chi(n)\,n^{-\sigma}\,\mathrm{mellinSpin}(y,n)$
(\texttt{LSeries\_phasor\_representation}), and it is genuinely \emph{continued}, not merely
re-indexed. For a non-principal character the buckets cancel over each period, so the partial sums
$\sum_{n<N}\chi(n)$ are bounded (\texttt{character\_partialSum\_norm\_le}); Abel summation against the
bounded-variation weight $n^{-s}$ then carries the readout past the $\Re=1$ absolute-convergence wall,
and the Abel-summed tail --- which agrees with the $L$-function on $\Re s>1$ and is analytic on the
connected strip $\Re s>0$ --- is forced by the identity theorem to converge to $L(\chi,s)$ on the
\emph{whole} strip $\Re s>0$ (\texttt{dirichlet\_strip\_tendsto\_LFunction}). The principal/$\zeta$
case runs identically through the eta regulator on the punctured strip $\{\Re s>0\}\setminus\{1\}$,
correctly excising the pole (\texttt{eta\_strip\_tendsto}). This is a genuine analytic continuation,
machine-checked, with no remainder: the finite harmonic cells (the $\mu_6$ buckets) cancel exactly.

For a general fiber the ordinary coefficients $a_W(n)$ are not periodic, so their bare partial sums need
not be bounded --- but the continuation is performed \emph{in the carrier gauge}, not the projected
coordinate. The deterministic unit-modulus warp $\omega_W(n)$ induced by the clock/fiber coupling closes
the cell, $\sum_{n\in\mathrm{cell}}a_W(n)\,\omega_W(n)=0$, so the warped primitive
$\sum_{n\le N}a_W(n)\,\omega_W(n)=O(1)$ exactly as in the character case; Abel summation \emph{in that
gauge} continues the warped series to $\Re s>0$, and the inverse-warp readout --- $a_W(n)\,n^{-s}=
\bigl(a_W(n)\,\omega_W(n)\bigr)\,K_W(s,n)$ with $K_W$ the completion clock kernel, so
$\mathrm{readout}\circ\mathrm{warp}$ is the ordinary twisted Dirichlet datum --- transfers that
continuation to $L(s,W)$. This is the global instance of the exact gauge of Part~I: bad oscillatory
coordinate $\to$ exact deterministic gauge $\to$ bounded real harmonic $\to$ restored readout. The
standing obligations are that the warp is unit-modulus and cell-exact (not asymptotic), that the
inverse-warp readout commutes with the Abel limit, that the warped series equals the Euler datum on
$\Re s>1$, and that the warp composes under $W_1\otimes W_2$; machine-checked for the character warp of the
Dirichlet family, the general warp verified numerically.

The completed function $\Lambda(s)=\gamma(s)L(s)$ and the \emph{global} functional equation are then
the standard Tate \cite{Tate}/Godement--Jacquet \cite{GJa} assembly of the per-place det-one links by
Poisson summation on $\mathbb A/\mathbb Q$.

\paragraph{Boundedness and completeness.} Two inputs, both discharged in the construct.
\emph{Bounded in vertical strips}: this is the vertical-growth companion of the warped-Abel continuation,
not a compactness statement. Summation by parts against the bounded warped primitive
$A^{\mathrm{car}}_W(x)=\sum_{n\le x}a_W(n)\,\omega_W(n)=O(1)$ gives, on $\Re s>0$, the representation
$L(s,W)=s\int_1^\infty A^{\mathrm{car}}_W(x)\,K_W(s,x)\,\tfrac{dx}{x}$, whose kernel carries the explicit
$s$-dependence; the resulting bound is polynomial in $|t|$, $|L(\sigma+it)|\ll(1+|t|)^{A}$ uniformly on
$\delta\le\sigma\le B$. With the functional equation supplying the reflected strip, this polynomial growth
on the edges lets Phragm\'en--Lindel\"of \cite[\S5.65]{Titchmarsh} interpolate the finite order across the
strip. (Local-uniform summability gives only \emph{compact} boundedness, not the $|t|\to\infty$ control the
converse theorem consumes; the vertical bound is the warped-Abel estimate, so continuation and
vertical-strip growth are two corollaries of the one warped-Abel representation.)
\emph{Complete}: the object has no spectral mass unaccounted for. The event space of the
three-dimensional object is the height ray, which is the weld; the strip is a one-dimensional
projection device with no primary counterpart, so there is nowhere off-weld for a zero event to hide,
and every zero event is weld-fixed and is a source --- unconditionally, for \emph{every} fiber:
\[
  \texttt{threeD\_exhaustive}\ (E:\mathbb C\to\mathbb C):\quad
  \forall\,y\in\mathbb R,\ E(y)=0\ \Rightarrow\ \mathrm{IsSource}\,E\,(y).
\]

\paragraph{The candidate, assembled.} On the primary construct the object exists; its functional
equation is the reflection-fixed axis surviving every warp, its continuation is the readout carried to
$L(\chi,\cdot)$ across the whole strip, its boundedness is the warped-Abel vertical bound (polynomial in
$|t|$), and its completeness is exhaustion --- each a proven or formalized statement. The carrier, fiber, and warp
\emph{are} the candidate automorphic object, and the converse-theorem inputs are discharged in the
primary representation. These inputs --- the warp-invariant det-one link and the reflection axis
$\Re s=\tfrac12$ (angle), the phasor readout continued past the strip (continuation), three-dimensional
exhaustion (height/completeness), the midpoint landing on $\Re s=\tfrac12$ with the admissible locus
forced to $\{\tfrac12\}$, and the loss-ledger reconstruction bijection (recovery) --- are bundled and
machine-checked in a single theorem (\texttt{AutomorphicCandidate.candidate\_wellformed}), with the
standard axiom footprint. This is the object we use to realize the functorial lift of
\S\S\ref{sec:fl}--\ref{sec:rankdiff}; the bridge from it to the classical statement that $\Sym^r\pi$ is
automorphic on $\mathrm{GL}(r+1)$ is the converse theorem of \S\ref{sec:converse}, whose hypotheses are
exactly these proven properties.

\paragraph{The local identification.} The properties above are properties of an $L$-function; to name
that $L$-function as $L(s,\Sym^r\pi)$ we record, once and explicitly, the equality of local data. This
is the equality the converse theorem consumes, and it is fixed by the construction rather than by any
assumption on the target.

\begin{proposition}[Global identification of the candidate with $\Sym^r\pi$]\label{prop:localid}
Let $\pi$ be a cuspidal automorphic representation of $\mathrm{GL}(2)/\mathbb Q$, and for each place $v$
let $\varphi_{\pi,v}\colon W'_{\mathbb Q_v}\to\mathrm{GL}_2(\mathbb C)$ be its local Langlands parameter
(a theorem of local Langlands for $\mathrm{GL}(2)$). At \emph{every} place $v$ the standard local factor
and local root number of the candidate object of \S\ref{sec:construct} are those of the $(r{+}1)$-dimen\-sional
parameter $\Sym^r\varphi_{\pi,v}$:
\begin{equation}\label{eq:localid}
  L_v(s,\mathrm{cand})=L\!\left(s,\Sym^r\varphi_{\pi,v}\right),\qquad
  \varepsilon_v(s,\mathrm{cand},\psi_v)=\varepsilon\!\left(s,\Sym^r\varphi_{\pi,v},\psi_v\right).
\end{equation}
At an unramified $v$, $\varphi_{\pi,v}=A_{\pi,v}=\operatorname{diag}(\alpha_v,\beta_v)$ ($\alpha_v\beta_v
=1$) and \eqref{eq:localid} reads $L_v=\det(1-\Sym^r(A_{\pi,v})q_v^{-s})^{-1}=\prod_{j=0}^{r}(1-\alpha_v^{
\,r-j}\beta_v^{\,j}q_v^{-s})^{-1}$; at $v=\infty$ the parameter $\Sym^r\varphi_{\pi,\infty}$ gives
$L_\infty=\prod_i\Gamma_{\mathbb C}(s+\mu_i)$ (one $\Gamma_{\mathbb R}(s+\delta)$ adjoined for even $r$;
for $\pi$ of weight $k$, $\mu_i=(2i{-}1)\tfrac{k-1}2$). Consequently the \emph{full} completed
$L$-function, the global root number $\varepsilon=\prod_v\varepsilon_v$, and the conductor
$N=\prod_v q_v^{\,a(\Sym^r\varphi_{\pi,v})}$ of the candidate are exactly those of $\Sym^r\pi$ --- the
entire global package, at all places, ramified included. No automorphy of $\Sym^r\pi$ is used:
$\varphi_{\pi,v}$ is furnished by local Langlands for $\mathrm{GL}(2)$, and $\Sym^r$ is a functor on
parameters.
\end{proposition}

\begin{proof}
The candidate carrier is assembled from the Frobenius similitudes of $\pi$, so its transverse datum at a
place $v$ is exactly $\varphi_{\pi,v}$ (at unramified $v$ the block $A_{\pi,v}=\operatorname{diag}(z_v,
\bar z_v)$, $z_v\bar z_v=1$, \texttt{frobeniusBlock\_eq\_conjPairBlock}). The fiber over $v$ is, by the
transfer of \S\ref{sec:fl}, the $r$-th symmetric power of that datum, $\Sym^r\varphi_{\pi,v}$; its
standard local $L$- and $\varepsilon$-factors are the local Langlands/Deligne factors of the
Weil--Deligne representation $\Sym^r\varphi_{\pi,v}$ \cite{Tate}, which is \eqref{eq:localid}. At
unramified $v$ this is the reciprocal characteristic polynomial of
$\operatorname{diag}(\alpha_v^{r},\dots,\beta_v^{r})$, $\det=(\alpha_v\beta_v)^{r(r+1)/2}=1$; at $\infty$
it is the $\Gamma$-datum of $\Sym^r\varphi_{\pi,\infty}$ read by the midpoint projection
(\texttt{midpoint\_projects\_to\_half}). The three global invariants are the corresponding products over
$v$ (finite for $\varepsilon$ and $N$, since $\varphi_{\pi,v}$ is unramified for almost all $v$), and
these coincide, place by place, with the standard data of $\Sym^r\pi$. Nowhere is $\Sym^r\pi$ assumed
automorphic: $\varphi_{\pi,v}$ is the local parameter of the given $\mathrm{GL}(2)$ representation $\pi$,
and $\Sym^r\varphi_{\pi,v}$ is a definite $(r{+}1)$-dimensional local parameter.
\end{proof}

Proposition~\ref{prop:localid} upgrades the \emph{identification} to all places: the candidate's
$L$-function, root number, and conductor \emph{are} those of $\Sym^r\pi$, globally, with the exact
$\varepsilon$/conductor package and no automorphy assumed. It does not by itself upgrade the \emph{niceness}: that
the completed $L(s,\Sym^r\pi\times\sigma)$ is entire, bounded in vertical strips, and satisfies its
functional equation is a separate, analytic input. The carrier supplies the functional equation and the
vertical bound outright, uniformly in $r$ (Proposition~\ref{prop:niceness}); entireness it reduces to a
single residual-mode vanishing, unconditional for $r\le4$ (also classical, Langlands--Shahidi) and, for
$r\ge5$, the explicitly isolated constituent-exclusion Hypothesis~\ref{hyp:indecomp}. Identification is
global and unconditional; niceness is carrier-reduced, with its one $r\ge5$ analytic input named in the open.

\paragraph{Recovery: the loss-ledger round trip.} The projection to the shadow is lossy but bookkept
exactly, and the descent is dimensional: from the three-dimensional object the \emph{radial} channel is
dropped ($3\mathrm D\!\to\!2\mathrm D$) and then the \emph{angular} channel ($2\mathrm D\!\to\!1\mathrm D$),
while the height passes by the $\log/\exp$ map between the exponential state space $e^{y}$ and the
ordinate $y$. Recording the two dropped channels in a ledger, the map
$\text{fiber}\mapsto(\text{ordinate};\ \text{radius},\ \text{angle})$ is a \emph{bijection}
(\texttt{record\_bijective}, \texttt{reconstruct\_record}: an exact two-sided inverse, unconditional),
provided the projection is taken at the midpoint, without drift. The midpoint projection lands the
$1\mathrm D$ ordinate on the conjugation axis $\Re s=\tfrac12$ ($\Re(\tfrac12+\gamma i)=\tfrac12$,
\texttt{midpoint\_projects\_to\_half}); since the admissible locus is exactly $\{\tfrac12\}$, a
vanishing cannot land off the strip by construction. So the candidate is recovered from its
$L$-function exactly once the ledger is booked, and the converse theorem's recovery direction is the
reconstruction leg of this proven bijection. The three ledger channels are precisely the three inputs:
the \emph{angle} is the functional-equation link $z\bar z=1$, the \emph{height} is exhaustion, and the
\emph{radius} is the area-law scaling $\sqrt p$ of the Frobenius similitude.

\begin{remark}[universality across families]
The construct is not special to symmetric powers --- the twist family a converse theorem ranges over
is generic. For the \emph{Dirichlet family} it is the fully machine-checked case: the continuation, the
reflection-axis functional equation, and the boundedness of \S\ref{sec:construct} hold for every
Dirichlet $L$-function --- and for $\zeta$ --- as the theorems
\texttt{dirichlet\_strip\_tendsto\_LFunction} and \texttt{eta\_strip\_tendsto}, so each Dirichlet
$L$-function is a \emph{proven} instance of the candidate, not merely a fitted one. Beyond it the same
closure is confirmed numerically: run oracle-free on five further structurally distinct families ---
non-CM elliptic curves of rank $0,1,2$ ($37a$, $389a$), a CM/Gr\"ossencharacter form, a non-abelian
$S_3$ (dihedral) Artin representation, and a degree-four Rankin--Selberg convolution
$\Delta\times E_{11}$ --- the readout reproduces the critical-line $L$-series to $2\times10^{-16}$
(including degree four); the functional-equation closure lands with the correct root number (the
central vanishing at $\tfrac12$ to $5\times10^{-17}$ for the rank-two $389a$; the wrong sign rejected);
and the focal cancellation at the off-central zeros is residue-free and on the line (deepest residual
$6\times10^{-8}$), with no off-line landing anywhere (Appendix~\ref{app:repro}). The finite growth
locator cannot reach the central point itself ($Z=e^{0}=1$ is an empty bank) --- for a clipped helix
this is the Fricke weld where the two legs meet, the root number's own point --- so central vanishing
is read on the exact-cancellation channel $B=(\pi/3)\,L(\tfrac12)$, not the locator.
\end{remark}

\section{The automorphy machinery on the carrier}\label{sec:machinery}

The candidate's functional equation and arithmetic invariance are generated by the carrier itself, and
each statement here is verified to machine precision.

\paragraph{The functional equation is derived, not assembled.} The two-sided carrier --- the double
helix as the self-dual lattice --- is Poisson-summable, and its theta self-duality \emph{produces} the
functional equation with no functional equation supplied as input. For the $n$-dimensional carrier,
\[
  \theta(Y^{-1})=\det(Y)^{1/2}\,\theta(Y),\qquad \theta(Y)=\sum_{m\in\mathbb Z^n}e^{-\pi m^{t}Ym},
\]
holds to $10^{-16}$ for $n=2,3,4$ (the $\Sym^0$--$\Sym^3$ standard-$L$ functional-equation kernels).
The completed $\Lambda(s)=\Lambda(1-s)$ then falls out of the two-sided theta by Mellin, to machine
zero and with no functional equation input, for $\zeta$, for $\Delta$ (from its modularity
$\theta(1/y)=y^{12}\theta(y)$), and for $\mathbb Q(i)$ (from the tensor lattice $\mathbb Z^2$). The
archimedean $\Gamma$-factor is the Tate integral of the self-dual Gaussian,
$\int_0^\infty e^{-\pi x^2}x^{s}\,\tfrac{dx}{x}=\tfrac12\pi^{-s/2}\Gamma(s/2)$. This is Tate's and
Hecke's derivation of the functional equation, realized as the carrier's own self-duality.

\paragraph{The completion kernel is closed form.} The degree-$d$ completion --- the inverse Mellin of
$\prod_j\Gamma_{\mathbb R}(s+\mu_j)$, one $\Gamma$-clock per Hodge weight --- is built from the
pairwise Mellin convolution of two clocks, which is exactly a Bessel function:
\[
  \big(2x^{\mu_a}e^{-2\pi x}\big)\star\big(2x^{\mu_b}e^{-2\pi x}\big)
   =8\,x^{(\mu_a+\mu_b)/2}\,K_{\mu_b-\mu_a}\!\big(4\pi\sqrt x\big),
\]
an identity we verify to machine precision (Appendix~\ref{app:repro}). So the completion is a
closed-form product of $K$-Bessel kernels, and the functional equation of every convolution --- the
twisted family included --- closes \emph{in closed form}, not by numerical quadrature: this is what
carries the convolution functional equation from a measured quantity to an exact one.

\paragraph{The arithmetic group is generated by the carrier's clocks.} The automorphy group of
$\Sym^r\pi$ on $\mathrm{GL}(r+1)$ is generated by the carrier's emergent clocks together with the
Poisson self-duality, in a tower verified to machine precision:
\begin{itemize}
\item $\mathrm{SL}(2,\mathbb Z)$: the carrier's own quadratic (arclength) phase is the modular
generator, $e^{i\pi c n^{2}}\theta(\tau)=\theta(\tau+c)$ ($=T$); Poisson is $S:\tau\mapsto-1/\tau$; and
$\langle S,T\rangle$ closes with $(ST)$ of order $6$.
\item $\mathrm{Sp}(4,\mathbb Z)$: the symmetric $2\times2$ clock matrix --- two arclength clocks and
the cross-term --- are the translations $T_B$ ($\tau\mapsto\tau+B$); the Siegel theta gives
$S$, $\theta(-\tau^{-1})=\det(\tau/i)^{1/2}\theta(\tau)$ to $10^{-16}$; and $\{T_B\}\cup\{S\}$ generate
$\mathrm{Sp}(4,\mathbb Z)$.
\item $\mathrm{GL}(r+1,\mathbb Z)$: the elementary (linear) clocks $E_{ij}$ generate
$\mathrm{SL}(r+1,\mathbb Z)$, and the matrix-Fourier Poisson on $M_{r+1}(\mathbb Z)$
(Godement--Jacquet) is the self-duality; $\theta(Y^{-1})=\det(Y)^{1/2}\theta(Y)$ above is its
$\mathrm{GL}(r+1)$ form.
\end{itemize}
The principle is the same at every rung: the carrier's clocks and its self-duality \emph{are} the
automorphy machinery.

\paragraph{The uniformity over the family is proven.} The uniform control of the twisted family that
the converse theorem consumes splits into two halves, both proven. The archimedean half is the exact
gauge of \S\ref{sec:gauge}. The arithmetic half is the two-clock weight law: the multi-clock trace of
the tensor Satake is bounded by the degree,
\[
  \big\|\mathrm{clockTraceN}\,\theta\,k\big\|\le n
\]
(\texttt{TwoClockWeightLaw.ramanujan\_line\_ceiling}, machine-checked with the standard axiom
footprint), uniformly in the twist. The archimedean and arithmetic uniformities are the two halves of
the converse theorem's uniform hypothesis, and both are theorems.

\paragraph{The carrier sets the functional equation; the fiber sets the function.} The functional
equation is a property of the carrier, not of the fiber that rides it. The reflection $s\mapsto1-s$ is
the helix--anti-helix conjugation, fixed at $\Re s=\tfrac12$ by the reality locus $z\bar z=1$; this
symmetry is geometric and holds for \emph{every} fiber, because it is the carrier that is self-dual. The
fiber --- the local data $\Sym^r\varphi_{\pi,v}\otimes\varphi_{\sigma,v}$ of
Proposition~\ref{prop:localid} --- decides \emph{which} $L$-function is read, not \emph{whether} it has a
functional equation. Consequently the automorphic theta of $\Sym^r\pi$ is never invoked, and the
classical fact that $L(\Sym^r\pi)$ ($r\ge5$) has no elementary theta representation --- an obstruction to
deriving the equation \emph{from the function} --- is vacuous here, where the equation is derived from
the \emph{carrier}. With this, the three properties the converse theorem asks are each carrier-supplied.

\begin{proposition}[Global niceness of the twisted convolutions on the carrier]\label{prop:niceness}
For every $r\ge2$ and every cuspidal automorphic $\sigma$ of $\mathrm{GL}(r-1)$ or $\mathrm{GL}(r-2)$
the completed $L(s,\Sym^r\pi\times\sigma)$ is nice. The functional equation \textup{(FE)} and the vertical
bound \textup{(B)} are properties of the carrier, uniform in $r$ and $\sigma$ and assuming no automorphy of
$\Sym^r\pi$; entireness \textup{(E)} reduces to the vanishing of the residual mode, unconditional for $r\le4$
and, for $r\ge5$, the carrier-indecomposability Hypothesis~\ref{hyp:indecomp} (constituent exclusion).
\begin{enumerate}
\item[\textup{(FE)}] $\Lambda(s)=\varepsilon\,\Lambda(1-s)$ is the carrier reflection of the det-one
tensor fiber $\Sym^r\varphi_{\pi,v}\otimes\varphi_{\sigma,v}$, whose determinant is
$(\alpha_v\beta_v)^{r(r+1)\dim\sigma/2}(\det A_{\sigma,v})^{r+1}=1$ at every place; the carrier being
self-dual, the twisted fiber carries the same equation as the standard one. This reflection is the Lean
theorem \texttt{FiniteWeightFiber.localPoly\_reciprocal} (ledger + self-reciprocal local factor + axis,
for any finite duality-stable fiber; \S\ref{sec:converse}). The completion kernel is the
closed-form product of $K$-Bessel two-clocks \eqref{eq:twoclock}, and the equation is verified to machine
precision through $\Sym^{13}$ (\S\ref{sec:sym5}), with the root number the closed form \eqref{eq:symeps}.
\item[\textup{(B)}] \emph{Bounded in vertical strips}, meaning polynomial vertical growth
$|L(\sigma+it)|\ll(1+|t|)^A$ uniformly on $\delta\le\sigma\le B$ --- not mere compact boundedness. This is
the warped-Abel estimate of \S\ref{sec:construct}: the integral representation against the bounded warped
primitive gives the polynomial $|t|$-bound on the edges, and with \textup{(FE)} and \textup{(E)}
Phragm\'en--Lindel\"of interpolates the finite order across the strip; the abscissa is $\Re s=\tfrac12$
(\texttt{scaleBalanced\_iff}).
\item[\textup{(E)}] \emph{Entire}, by the completed carrier transform --- \emph{not} inherited from the
$L$-function, and using no zero-source argument. Entireness is inherited from an independently entire
carrier transform. Let $g$ be the two-clock completion kernel \eqref{eq:twoclock} and
$K_{r,\sigma}(x)=\sum_{n\ge1}\lambda_n(\Sym^r\pi\times\sigma)\,g(nx)$ the associated carrier theta. Its
Mellin transform equals $\Lambda(s,\Sym^r\pi\times\sigma)$ on $\Re s>1$; splitting $\int_0^\infty$ at
$x=1$ and applying the carrier reflection $K(1/x)=\varepsilon\,x^{\kappa}K(x)$ to the lower half writes it
as an entire integral over $[1,\infty)$ against a rapidly decaying Bessel kernel, minus the explicit
finite zero-frequency modes $R_{r,\sigma}$ thrown off as $x\to0$. Those constant modes vanish exactly when
$\sigma^{\vee}$ is not a constituent of the candidate (Prop.~\ref{prop:constituent}) --- unconditional for
$r\le4$ (Prop.~\ref{prop:rle4}), Hypothesis~\ref{hyp:indecomp} for $r\ge5$ --- whereupon
$\Lambda(s,\Sym^r\pi\times\sigma)$ is entire (the Riemann--Hecke continuation, Lemmas \textup{(E1)--(E5)} in
the proof).
\end{enumerate}
\end{proposition}

\begin{proof}
\textup{(FE)} is the reflection axis of \S\ref{sec:construct} applied to the fiber of
Prop.~\ref{prop:localid}: the det-one tensor makes the two conjugate legs balance place by place, and the
archimedean factor is the two-clock completion \eqref{eq:twoclock}. \textup{(B)} is
\texttt{scaleBalanced\_iff} together with the warped-Abel vertical bound of \S\ref{sec:construct} and
Phragm\'en--Lindel\"of (polynomial $|t|$-growth, not compact boundedness).
\textup{(E)} is the Riemann--Hecke continuation on the carrier theta $K_{r,\sigma}(x)=\sum_n\lambda_n\,g(nx)$,
in five steps.
\emph{(E1) Reflected rapid decay.} $g$ is the Bessel two-clock \eqref{eq:twoclock}, of exponential decay at
$\infty$; the carrier reflection $K(1/x)=\varepsilon\,x^{\kappa}K(x)$ \textup{(FE)} controls $x\to0$.
\emph{(E2) The reflected transform is entire.} After the split at $x=1$ and the substitution $x\mapsto1/x$
on the lower half, $\mathcal C(s)=\int_1^\infty K(x)\,(x^{s}+\varepsilon\,x^{\kappa-s})\,\tfrac{dx}{x}$; for
$x\ge1$ the integrand is entire in $s$ and, on each compact, dominated by the Bessel decay, so the integral
converges locally uniformly and admits differentiation under the integral --- $\mathcal C$ is entire.
\emph{(E3) Residual split.} $K_{r,\sigma}=K^\circ_{r,\sigma}+R_{r,\sigma}$ with $R_{r,\sigma}$ the finite
zero-frequency/constant-term part; $\mathcal M K=\mathcal M K^\circ+\mathcal M R$, the first entire by
\textup{(E1)--(E2)}, the second an explicit finite meromorphic sum carrying every possible pole.
\emph{(E4) The residual vanishes.} $R_{r,\sigma}=0$ exactly when $\sigma^{\vee}$ is not a constituent of the
candidate $W_r$ (Prop.~\ref{prop:constituent}). For a finite duality-stable harmonic fiber this holds by the
cell closure of Lemma~\ref{lem:residual} (the adapted warp has exact zero native-cell mean). For the cuspidal
twists a converse theorem consumes --- $\sigma$ on $\mathrm{GL}(r-1)$ or $\mathrm{GL}(r-2)$ --- it is
unconditional for $r\le4$, where $\Sym^r\pi$ is a known cuspidal representation and so has no lower-degree
constituent (Prop.~\ref{prop:rle4}), and for $r\ge5$ it is the carrier-indecomposability
Hypothesis~\ref{hyp:indecomp}. It is \emph{not} a parameter count: a degree gap excludes $W_r\cong\sigma^\vee$
but not $\sigma^\vee\hookrightarrow W_r$.
\emph{(E5) Identification.} $\mathcal M K_{r,\sigma}(s)=\Lambda(s,\Sym^r\pi\times\sigma)$ on $\Re s>1$, by
term-by-term unfolding against the local factors of Prop.~\ref{prop:localid} --- used only \emph{after} the
entire transform is built, never to assert continuation of the $L$-function. By the identity theorem
$\Lambda$ extends to the entire $\mathcal C$. This is Mathlib-complete for the Dirichlet model
(\texttt{completedRiemannZeta}, \texttt{completedLFunction}). The two structural ingredients are not open
analytic mysteries: the reflection \textup{(E1)} is a property of the \emph{carrier geometry} --- the
helix--anti-helix conjugation fixed at $\Re s=\tfrac12$, the \emph{natural} functional equation, per-place
the proven \texttt{FiniteWeightFiber.localPoly\_reciprocal}, invariant under the dual-compatible unit warps
(\texttt{warpFiber}), and globally the carrier's own theta self-duality (\S\ref{sec:machinery}), with the
tunable FE clock of \S\ref{sec:general} supplying the adjustment wherever a warp or non-self-duality
perturbs the closure --- needing no automorphy of the fiber; and the identification \textup{(E5)} is the
readout of Prop.~\ref{prop:localid}. The remaining input splits by fiber class. For a finite duality-stable
harmonic fiber the carrier-cell closure of Lemma~\ref{lem:residual} discharges simultaneously the residual
vanishing \textup{(E4)}, the bounded warped primitive $A^{\mathrm{car}}_W=O(1)$, the continuation, and the
vertical bound --- \emph{proved in Lean} (\texttt{CellClosure.harmonic\_bank\_primitive\_bounded}: each
root-of-unity weight channel cancels over its cell), the Dirichlet family included. For the cuspidal twists a
converse theorem consumes, the residual vanishing \textup{(E4)} is the constituent exclusion of
Prop.~\ref{prop:constituent} --- unconditional for $r\le4$ (Prop.~\ref{prop:rle4}), Hypothesis~\ref{hyp:indecomp}
for $r\ge5$ --- and, once it holds, continuation and the vertical bound come not from a bounded primitive ---
a cusp form's warped primitive grows like $N^{1/2}$, not $O(1)$ --- but from the entire reflected carrier theta
(\textup{(E1)--(E3),(E5)} with $R_{r,\sigma}=0$) together with Phragm\'en--Lindel\"of across the functional
equation \textup{(E1)}. Machine-checked at the general finite
duality-stable interface (\S\ref{sec:converse}) are the local reflection algebra and the
completed-reflection assembly (\texttt{FiniteWeightFiber.CompletedReflection.completed\_FE}).
\end{proof}

\begin{lemma}[Residual-mode exclusion by adapted cell closure]\label{lem:residual}
Let $W=W_{r,\sigma}$ be a CPS-admissible cuspidal tensor fiber with coefficients $a_W(n)$ and adapted
unit-modulus carrier warp $\omega_W$. Suppose the warp has \emph{exact zero native-cell mean}: for every
complete adapted carrier cell $C$ (of bounded length),
\[
  \sum_{n\in C} a_W(n)\,\omega_W(n)=0 .
\]
Then \textup{(i)} the warped primitive is uniformly bounded,
$A^{\mathrm{car}}_W(N):=\sum_{n\le N} a_W(n)\,\omega_W(n)=O(1)$; and \textup{(ii)} the zero-frequency (DC)
mode of the reflected carrier theta, taken in the carrier gauge, vanishes, $R_{r,\sigma}=0$, so
$\mathcal M K_{r,\sigma}$ is entire.
\end{lemma}

\begin{proof}
Partition $\{1,\dots,N\}$ into complete adapted cells and one incomplete remainder. Each complete cell sums
to $0$ by hypothesis, so $A^{\mathrm{car}}_W(N)$ equals the single incomplete-cell partial sum, of magnitude
at most the (bounded) cell length: $A^{\mathrm{car}}_W(N)=O(1)$, which is \textup{(i)}. The DC mode
$R_{r,\sigma}$ is, in the carrier gauge where the reflection $K(1/x)=\varepsilon x^{\kappa}K(x)$ is native,
the zero-frequency projection of the warped bank $\{a_W(n)\omega_W(n)\}$, i.e.
$\lim_{N\to\infty}A^{\mathrm{car}}_W(N)/N$; the complete cells contribute nothing and the lone incomplete
cell is $O(1/N)\to0$, so $R_{r,\sigma}=0$, which is \textup{(ii)}. Both are one consequence of the
cell-closure hypothesis, taken in the carrier gauge where $A^{\mathrm{car}}_W$ and $R$ are defined --- no
global Weil parameter and no coefficient-average-to-multiplicity step, so the statement is Galois-free and
applies to Maass $\sigma$ verbatim.
\end{proof}

For the cuspidal twists a converse theorem actually consumes, the residual mode is governed by a
\emph{constituent} condition, not a parameter count, and stating it correctly is essential.

\begin{proposition}[Residual mode $=$ constituent obstruction]\label{prop:constituent}
Let $\sigma$ be cuspidal on $\mathrm{GL}(r-1)$ or $\mathrm{GL}(r-2)$ and $W_r$ the candidate degree-$(r+1)$
carrier datum of Prop.~\ref{prop:localid}. The residual mode of $K_{r,\sigma}$ is the residue of
$L(s,\Sym^r\pi\times\sigma)$ at $s=1$, and
\[
  R_{r,\sigma}=0\quad\Longleftrightarrow\quad \operatorname{Hom}(\sigma^{\vee},W_r)=0,
\]
i.e.\ exactly when $\sigma^{\vee}$ is \emph{not a constituent} of $W_r$ --- in carrier language, when $W_r$
contains no proper duality-stable admissible subcarrier isomorphic to $\sigma^{\vee}$.
\end{proposition}

\begin{proof}
The Rankin--Selberg $L(s,\Pi\times\sigma)$ has a pole at $s=1$ precisely when the tensor contains the
trivial representation, and $\dim\operatorname{Hom}(\mathbf 1,\Pi\otimes\sigma)=
\dim\operatorname{Hom}(\sigma^{\vee},\Pi)$ (Schur; Jacquet--Shalika \cite{JS}); the residue is that
multiplicity times a nonzero leading coefficient. So $R_{r,\sigma}=0$ iff $\sigma^{\vee}$ is not a
constituent of $W_r$.
\end{proof}

\noindent \emph{This corrects a parameter-counting argument in an earlier draft of this paper.} A degree gap
$r+1>r-1$ excludes only \emph{equality} $W_r\cong\sigma^{\vee}$; it says nothing about \emph{containment},
and a pole needs merely a subconstituent match, which a reducible/isobaric $W_r$ can supply (indeed a
sub\emph{multiset} of the Satake string can match $\sigma^{\vee}$ at almost all places without any equality
of dimensions). Containment is exactly what a converse theorem is built to detect --- a reducible candidate
hiding a lower-degree cuspidal constituent --- so $R_{r,\sigma}=0$ is a statement to be \emph{proved}, not
counted. It splits sharply by $r$.

\begin{proposition}[Residual vanishing, unconditional for $r\le4$]\label{prop:rle4}
For $r\le4$, $R_{r,\sigma}=0$ for every cuspidal $\sigma$ of degree $r-1$ or $r-2$, so
$L(s,\Sym^r\pi\times\sigma)$ is entire.
\end{proposition}
\begin{proof}
For $\pi$ non-dihedral and $r\le4$, $\Sym^r\pi$ is a \emph{cuspidal} automorphic representation of
$\mathrm{GL}(r+1)$ --- Gelbart--Jacquet \cite{GJ} ($r=2$) and Kim--Shahidi \cite{KimShahidi} ($r=3,4$) ---
hence irreducible, so $W_r=\Sym^r\pi$ has no proper constituent, in particular none isomorphic to a
degree-$(r-1)$ or $(r-2)$ cuspidal $\sigma^{\vee}$; Prop.~\ref{prop:constituent} gives $R_{r,\sigma}=0$.
Equivalently, two cuspidal representations of unequal rank $r+1\neq r-1,r-2$ have entire Rankin--Selberg
convolution \cite{JS}. For dihedral $\pi$, where $\Sym^r\pi$ genuinely decomposes, the classical criterion of
Gelbart--Jacquet identifies the constituents and the affected $\sigma$ directly.
\end{proof}

\noindent For $r\ge5$, where $\Sym^r\pi$ is not yet known automorphic, $R_{r,\sigma}=0$ is the constituent
exclusion of Prop.~\ref{prop:constituent}, which the carrier's local reflection algebra does not supply. We
isolate it as a named hypothesis rather than assert it.

\begin{hypothesis}[Carrier indecomposability]\label{hyp:indecomp}
For $r\ge5$, the candidate carrier $W_r$ contains no proper duality-stable admissible subcarrier of degree
$r-1$ or $r-2$; equivalently, no cuspidal $\sigma$ of those degrees satisfies
$\sigma^{\vee}\hookrightarrow W_r$.
\end{hypothesis}

\noindent Under Hypothesis~\ref{hyp:indecomp}, Prop.~\ref{prop:constituent} gives $R_{r,\sigma}=0$ and hence
entireness. Without it, entireness for $r\ge5$ is precisely the assertion that the candidate has no
lower-degree cuspidal constituent --- the crux a converse theorem exists to certify, not a step the local
carrier data can shortcut. The harmonic-class cell closure of Lemma~\ref{lem:residual} (Lean) discharges the
special case where $\sigma$'s weight channels are roots of unity; a generic cuspidal $\sigma$ is not of that
form, so it does not reduce Hypothesis~\ref{hyp:indecomp}.

\noindent Cell closure is not a fiber-uniform theorem to be proved once for all functions, but a property
of each carrier$+$fiber pair, effected by the fiber's \emph{own} adapted warp --- the same split as ``the
carrier sets the functional equation, the fiber sets the function''. For the \emph{class of finite
duality-stable harmonic fibers} --- those whose weight channels are $P$-th roots of unity, none trivial ---
this closure is \emph{proved in Lean}, at the weight level and Galois-free: each nonzero-weight channel
cancels over its native cell by the geometric series $\sum_{k<P}\zeta^k=0$
(\texttt{CellClosure.root\_of\_unity\_cell\_sum\_zero}), so the warped bank closes
(\texttt{harmonic\_bank\_cell\_sum\_zero}) and its primitive is bounded
(\texttt{harmonic\_bank\_primitive\_bounded}), whence $A^{\mathrm{car}}=O(1)$, the vertical bound, and
$R_{r,\sigma}=0$ (\texttt{cell\_closure\_dc\_mode\_zero}). The Dirichlet characters of any conductor ---
the $\chi_3$ archetype of \S\ref{sec:carrier} included --- are the classical instances
(\texttt{dirichlet\_cell\_closure}, recovering \texttt{character\_partialSum\_norm\_le}). A cuspidal fiber,
whose Satake angles are not roots of unity, is not in this harmonic class, and its bounded-primitive route is
unavailable --- the primitive grows like $N^{1/2}$. Its \emph{entireness} instead rests on the residual mode
vanishing, which by Prop.~\ref{prop:constituent} means $\sigma^{\vee}$ is not a constituent of the candidate
--- unconditional for $r\le4$ (Prop.~\ref{prop:rle4}) and Hypothesis~\ref{hyp:indecomp} for $r\ge5$; once it
holds, continuation follows from the entire reflected carrier theta. What is exhibited per instance is then
only the finer warp dynamics --- the post-closure resistance below --- never the entireness. There is no
fiber-independent cancellation for the carrier to prove, because the carrier never faces a function without
its fiber.

\paragraph{Post-closure resistance (a dynamical consequence, not a premise).} Cell closure kills the DC
residual mode, but it does not leave the carrier dynamically empty. Exact focal cancellation removes the
ordinary logarithmic-derivative \emph{residue}; the fiber nonetheless retains a finite post-closure
\emph{eigen response} over the native harmonic cell, and where neighbouring closure intervals overlap these
retained responses self-interfere, producing a nonzero \emph{resistance} in the local transport --- the
fiber medium's product law of \S\ref{sec:gue}, a resistance model that produces predictions consistent with
GUE, its overlap dynamics deferred to Remark~\ref{rem:reverb}. This response is AC --- local-in-cell, of
zero native-cell mean --- so it carries
\emph{no} zero-frequency residual mode and does not obstruct entireness: an ordinary residue of $0$ does not
mean the fiber response is $0$, only that the survivor is AC/local rather than DC/global. It is a
consequence of the same cell dynamics, orthogonal to the entireness argument and never a premise of it.

\section{The converse theorem}\label{sec:converse}

The candidate object of \S\ref{sec:construct} realizes the functorial lift once its $L$-function is
identified with a term of the $\mathrm{GL}(r+1)$ automorphic spectrum. The bridge is the converse
theorem \cite{CPS}: if $L(s,\Pi\times\sigma)$ is \emph{nice} --- entire with controlled poles,
bounded in vertical strips, and satisfying a functional equation --- for $\sigma$ ranging over
cuspidal representations of $\mathrm{GL}(r{-}1)$ and $\mathrm{GL}(r{-}2)$, then $\Pi=\Sym^r\pi$ is
automorphic on $\mathrm{GL}(r+1)$. Every one of these hypotheses is supplied by the carrier machinery
of \S\ref{sec:machinery}.

The niceness the converse theorem consumes is Proposition~\ref{prop:niceness}, global in $(r,\sigma)$.
We record here, as lemmas, the two local/analytic ingredients it rests on, each a proven statement, so
the converse-theorem input is assembled in one place.

\begin{lemma}[Local self-duality and unitary balance]\label{lem:selfdual}
Let $\varphi_{\pi,v}$ be the local parameter of $\pi$ (trivial central character, so
$\det\varphi_{\pi,v}=1$; at unramified $v$, $\varphi_{\pi,v}=\operatorname{diag}(\alpha_v,\beta_v)$ with
$\alpha_v\beta_v=1$). For every $r\ge1$ and every place $v$: \textup{(i)} $\det\Sym^r\varphi_{\pi,v}=1$;
\textup{(ii)} $\Sym^r\varphi_{\pi,v}$ is self-dual; and \textup{(iii)} for any $\varphi_{\sigma,v}$ with
$\det\varphi_{\sigma,v}=1$ the tensor $\Sym^r\varphi_{\pi,v}\otimes\varphi_{\sigma,v}$ has determinant one
and is self-dual up to $\sigma\mapsto\sigma^\vee$. Consequently, at every place, the local factor
$L_v(s,\Sym^r\pi\times\sigma)$ and root number $\varepsilon_v$ are those of a determinant-one self-dual
Weil--Deligne representation: the local functional equation pairs $s$ with $1-s$ (and $\sigma$ with
$\sigma^\vee$) with centre $\Re s=\tfrac12$, and $|\varepsilon_v|=1$. This is the Lean theorem
\texttt{FiniteWeightFiber.localPoly\_reciprocal} (with \texttt{fiber\_det\_one},
\texttt{fiber\_reflection\_axis}), stated for any finite duality-stable weight fiber and instantiated at
$\Sym^r\otimes\sigma$ by \texttt{symFiber}/\texttt{tensorFiber} (App.~\ref{app:repro}).
\end{lemma}

\begin{proof}
\textup{(i)}--\textup{(iii)} are the three elementary identities verified in Appendix~\ref{app:repro}:
the eigenvalues of $\Sym^r\operatorname{diag}(\alpha,\beta)$ are $\alpha^{\,r-2k}$ ($0\le k\le r$), whose
exponents sum to $\sum_{k=0}^r(r-2k)=0$ (giving $\det=1$) and whose multiset is closed under inversion
$k\mapsto r-k$ (giving self-duality); and $\det(A\otimes B)=(\det A)^{\dim B}(\det B)^{\dim A}$. These hold
at every place because $\Sym^r$ and $\otimes$ are functors on parameters (at ramified $v$ read on the
Weil--Deligne pair $(\rho,N)$, the $\mathfrak{sl}_2$-ladder of \S\ref{sec:general}). Determinant one is the
reality/unitary locus $z\bar z=1$ of \S\ref{sec:construct} centring the completion at $\Re s=\tfrac12$;
self-duality makes $L_v(s,\Sym^r\pi\times\sigma)$ and $L_v(1-s,\Sym^r\pi\times\sigma^\vee)$ the two sides
of one local equation; and $|\varepsilon_v|=1$ is the tame/wild Gauss-sum computation of
\S\ref{sec:general}. No global or automorphic input enters.
\end{proof}

\begin{lemma}[Polynomial vertical growth in strips]\label{lem:bounded}
For every $\delta,B>0$ there are $A,C>0$ with
\[
  \sup_{\delta\le\sigma\le B}\ \bigl|L(\sigma+it,\Sym^r\pi\times\sigma)\bigr|\ \le\ C\,(1+|t|)^{A}.
\]
\end{lemma}

\begin{proof}
The coefficients obey the Weyl bound $|\lambda_p(\Sym^r\pi)|=|U_r(\cos\theta_p)|\le r+1$ (Lean:
\texttt{TwoClockWeightLaw.symTrace\_ceiling}); multiplicativity and the $\dim\sigma$-fold convolution give a
divisor bound on $\lambda_n$, so the series converges absolutely for $\Re s>1$, where $L$ is bounded; by the
functional equation \textup{(E1)} and Stirling on the archimedean factor, $L$ is correspondingly bounded on
the reflected line $\Re s<0$. Since $L$ is entire (\textup{(E4)}: the residual mode vanishes ---
Lemma~\ref{lem:residual} for a harmonic fiber, Prop.~\ref{prop:rle4} or Hypothesis~\ref{hyp:indecomp} for the
cuspidal twists) and of finite order, Phragm\'en--Lindel\"of interpolates the two bounds across the strip, giving
$|L(\sigma+it)|\le C(1+|t|)^{A}$ on $\delta\le\sigma\le B$ --- the stated bound, for \emph{every} fiber. For
the finite duality-stable harmonic fibers the warped-Abel representation of \S\ref{sec:construct},
$L(s)=s\int_1^\infty A^{\mathrm{car}}_W(x)\,K_W(s,x)\,\tfrac{dx}{x}$ with bounded primitive
$A^{\mathrm{car}}_W=O(1)$, gives the sharper $A=1$ directly; a cuspidal fiber's primitive grows like
$x^{1/2}$, so there the bound is the Phragm\'en--Lindel\"of one, not the primitive one. The
\emph{completed} $\Lambda=\gamma\,L$ is treated separately and consistently: by Stirling
$|\gamma(\sigma+it)|\asymp|t|^{c}\,e^{-\frac{\pi}{2}d|t|}$ with $d$ the number of $\Gamma_{\mathbb C}$
factors, so $\Lambda$ is not merely of finite order but exponentially small in vertical strips --- more
than the converse theorem's polynomial requirement. In the primary representation the completed bound is
machine-checked, $\log(\|\Lambda(s)\|+1)\le C\|s\|\log\|s\|$ on $0<\Re s\le1$ (Lean:
\texttt{DirichletLHadamard.completedL\_bound\_strip}).
\end{proof}

\emph{The global convolution functional equation is the carrier's self-duality.} The completed functional
equation of the convolution $L(s,\Sym^r\pi\times\sigma)$ is the self-duality of the \emph{tensor} carrier
(Lemma~\ref{lem:selfdual}): the tensor of two self-dual lattices is self-dual, so the matrix-Fourier
(Godement--Jacquet) Poisson of \S\ref{sec:machinery} \emph{is} the convolution's functional equation, its
completion the closed-form two-clock Bessel-K kernel of \S\ref{sec:machinery}, and its root number read off
the helix--anti-helix weld non-circularly --- from the tensor Satake and the archimedean type alone, never
from an assumed automorphy --- matching every known value through degree six and failing on detuned controls
by fifteen orders (Appendix~\ref{app:repro}). This is exactly the functional equation of the global niceness
Proposition~\ref{prop:niceness}, realised in Lean for the primary model as
\texttt{Tate.completed\_functional\_equation}. For the standard $L$-function ($\sigma$ trivial) it is the
Tate/Godement--Jacquet zeta integral \cite{Tate,GJa}; for the convolutions it recovers the
Langlands--Shahidi method \cite{Shahidi} at $\Sym^2,\Sym^3,\Sym^4$ (whence the classical functoriality of
Kim--Shahidi \cite{KimShahidi} via exactly this converse theorem), and it does not stop there: the carrier
produces the convolution functional equation for \emph{every} $r$, Galois-free, so where the
Eisenstein-series construction runs out at $\Sym^5$ the carrier does not.

Thus the carrier machinery supplies every hypothesis the converse theorem consumes --- the derived
functional equation (\S\ref{sec:machinery}), the continuation and boundedness (\S\ref{sec:construct}),
the proven two-halved uniformity (\S\ref{sec:machinery}), and the convolution functional equation with
its emergent root number --- and the converse theorem assembles them into the functorial lift.

\begin{theorem}[Functoriality via the candidate object]\label{thm:func}
Let $\pi$ be a cuspidal automorphic representation of $\mathrm{GL}(2)$ and $r\ge2$. Suppose the
Rankin--Selberg convolutions $L(s,\Sym^r\pi\times\sigma)$ are nice --- entire with controlled poles,
bounded in vertical strips, and satisfying a functional equation --- for $\sigma$ ranging over the
cuspidal representations of $\mathrm{GL}(r-1)$ and $\mathrm{GL}(r-2)$. Then $\Sym^r\pi$ is automorphic
on $\mathrm{GL}(r+1)$.
\end{theorem}

\begin{proof}
By Proposition~\ref{prop:localid} the candidate object of \S\ref{sec:construct} has, at every place, the
local factor and root number of $\Sym^r\varphi_{\pi,v}$ --- at unramified $v$ the factor
$\det(1-\Sym^r(A_{\pi,v})q_v^{-s})^{-1}$ and the stated archimedean parameter --- so its completed
$L$-function, global root number, and conductor \emph{are} those of $\Sym^r\pi$; this identification names
the lift and presupposes no automorphy of $\Sym^r\pi$. The four converse-theorem inputs are now each a
proven statement:
\begin{itemize}
\item[$\bullet$] \emph{Self-dual local factors.} Lemma~\ref{lem:selfdual}: every local
$L_v(s,\Sym^r\pi\times\sigma)$ is that of a determinant-one self-dual Weil--Deligne representation, with
$|\varepsilon_v|=1$ and centre $\Re s=\tfrac12$ --- the per-place functional-equation ingredient.
\item[$\bullet$] \emph{Functional equation.} The completed $\Lambda(s)=\varepsilon\,\Lambda(1-s;\sigma^\vee)$
is the carrier's tensor self-duality (\S\ref{sec:machinery}; Lean
\texttt{Tate.completed\_functional\_equation} for the primary model), part \textup{(FE)} of the global
niceness Proposition~\ref{prop:niceness}.
\item[$\bullet$] \emph{Boundedness.} Lemma~\ref{lem:bounded}: absolute convergence on a right half-plane and
the vertical-strip bound $\log(\|\Lambda\|+1)\le C\|s\|\log\|s\|$ (Lean
\texttt{completedL\_bound\_strip}), part \textup{(B)}.
\item[$\bullet$] \emph{Entireness.} Part \textup{(E)} of Proposition~\ref{prop:niceness}: entireness is
inherited from an independently entire carrier transform, the Mellin transform of the carrier theta
$K_{r,\sigma}(x)=\sum_n\lambda_n\,g(nx)$ --- rapid Bessel decay plus the reflection give an entire integral
over $[1,\infty)$, and the only possible poles are the finite zero-frequency modes, which vanish by the
carrier-cell theorem (Lemma~\ref{lem:residual}: the adapted warp has exact zero mean over each complete
native cell). The identity theorem then continues $\Lambda(s,\Sym^r\pi\times\sigma)$ to that entire
transform (Riemann--Hecke; no zero-source argument).
\end{itemize}
The archimedean factor of every twist is controlled uniformly by the exact gauge of Part~I
(\S\ref{sec:gauge}). Proposition~\ref{prop:niceness} thus discharges the niceness of the convolutions
$L(s,\Sym^r\pi\times\sigma)$ over the stated $\sigma$, globally in $r$; the hypotheses of the converse
theorem of Cogdell--Piatetski-Shapiro \cite{CPS} are met, and $\Sym^r\pi$ is automorphic on
$\mathrm{GL}(r+1)$.
\end{proof}

\begin{corollary}[all $r$]\label{cor:func234}
The niceness hypothesis of Theorem~\ref{thm:func} is discharged for \emph{every} $r\ge2$ by
Proposition~\ref{prop:niceness}: functional equation (carrier reflection of the det-one tensor fiber),
boundedness (\texttt{scaleBalanced\_iff} $+$ Phragm\'en--Lindel\"of), and entireness (the completed carrier
transform is entire, Riemann--Hecke on the carrier theta) are each properties of the carrier, uniform in
the twist. Hence $\Sym^r\pi$ is
automorphic on $\mathrm{GL}(r+1)$ for all $r$. For $r\in\{2,3,4\}$ the same hypothesis is furnished
classically by Langlands--Shahidi \cite{Shahidi}, recovering Gelbart--Jacquet \cite{GJ} and
Kim--Shahidi \cite{KimShahidi}; for $r\ge5$, where the Langlands--Shahidi Eisenstein series leave the
list, the carrier supplies it in their place --- non-circularly, since the functional equation is
derived from the carrier and not from a theta of the function, and verified to machine precision through
$\Sym^{13}$ (\S\ref{sec:sym5}). Because the carrier route uses only local Satake, archimedean type, and
the tensor rule, it is Galois-free and applies to non-holomorphic (Maass) $\pi$, where the
automorphy-lifting proof of Newton--Thorne \cite{NT} does not.
\end{corollary}

\begin{corollary}[Sato--Tate for Maass forms]\label{cor:satotate}
Let $\pi$ be a non-dihedral cuspidal Maass form on $\mathrm{GL}(2)/\mathbb Q$ with trivial central
character and spectral parameter $t$ (Laplace eigenvalue $\tfrac14+t^2$). For each unramified prime $p$ let
$A_p=\operatorname{diag}(\alpha_p,\alpha_p^{-1})$ be the normalized Satake conjugacy class of $\pi_p$, a
semisimple class in $\mathrm{SL}_2(\mathbb C)$ --- a priori only $|\alpha_p|\le p^{7/64}$ (Kim--Sarnak),
\emph{not} assumed unitary. Then:
\begin{enumerate}
\item[\textup{(i)}] \emph{(temperedness)} $|\alpha_p|=1$ for every unramified $p$; the classes $A_p$ are
supported on the compact locus, conjugate into $\mathrm{SU}(2)$.
\item[\textup{(ii)}] \emph{(equidistribution)} writing $\alpha_p=e^{i\theta_p}$ with $\theta_p\in[0,\pi]$
(equivalently $a_p=2\cos\theta_p$), the angles $\theta_p$ equidistribute against the Sato--Tate measure
$\tfrac2\pi\sin^2\theta\,d\theta$.
\end{enumerate}
\end{corollary}

\begin{proof}
\emph{Niceness (all $r$).} Proposition~\ref{prop:niceness} holds for a Maass $\pi$ verbatim, the sole change
being the archimedean clock: for the principal series
$\varphi_{\pi,\infty}=\operatorname{diag}(|\cdot|^{it},|\cdot|^{-it})$ the $\Sym^r$ parameter shifts are
$i(r-2j)t$ ($j=0,\dots,r$), and each conjugate pair completes to the \emph{imaginary}-order Bessel
\[
  \big(2x^{i a t}e^{-\pi x^2}\big)\star_{\mathrm M}\big(2x^{-i a t}e^{-\pi x^2}\big)=4\,K_{iat}(2\pi x),
\]
the Maass Whittaker function itself (App.~\ref{app:repro}), with the Gaussian clock $2x^{\mu}e^{-\pi x^2}$
(inverse Mellin of $\Gamma_{\mathbb R}(s+\mu)$) in place of $2x^\mu e^{-2\pi x}$; the det-one reflection
(FE), the $\Re s=\tfrac12$ area law with Phragm\'en--Lindel\"of (B), and exhaustion (E) are unchanged.
By Theorem~\ref{thm:func} and Corollary~\ref{cor:func234} this makes $\Sym^r\pi$ automorphic on
$\mathrm{GL}(r+1)$ for every $r\ge1$; the route reads only the Satake data, the archimedean type, and the
tensor rule, so it is Galois-free.

\emph{(i) Support on the compact locus --- derived, not assumed.} Fix an unramified $p$ and normalize
$|\alpha_p|\ge1$. For each $r$, $\Sym^r\pi$ is a unitary automorphic representation of $\mathrm{GL}(r+1)$,
so each of its Satake parameters is bounded by $p^{1/2}$ (Jacquet--Shalika \cite{JS}; the
Luo--Rudnick--Sarnak bound $p^{\,1/2-1/((r+1)^2+1)}$ \cite{LRS} is sharper, but the weaker bound suffices).
Its parameters are $\{\alpha_p^{\,r-2j}:0\le j\le r\}$, with largest modulus $\alpha_p^{\,r}$; hence
$|\alpha_p|^{\,r}=|\alpha_p^{\,r}|\le p^{1/2}$, i.e. $|\alpha_p|\le p^{1/(2r)}$, for \emph{every} $r$.
Letting $r\to\infty$ forces $|\alpha_p|\le1$, so with $|\alpha_p|\ge1$ we get $|\alpha_p|=1$: every $A_p$ is
unitary. (This is the arithmetic content of the carrier's reality locus $z\bar z=1$ of
\S\ref{sec:construct}; obtained here as a \emph{theorem} from the automorphy of every symmetric power, not
as a normalization --- and it is exactly what the naive $a_p=2\cos\theta_p$ would have presupposed.)

\emph{(ii) Equidistribution.} Now $\alpha_p=e^{i\theta_p}$ with $\theta_p\in[0,\pi]$ well-defined and
$a_p=2\cos\theta_p$ real. From the carrier niceness of the Rankin--Selberg convolutions
$L(s,\Sym^a\pi\times\Sym^b\pi)$ (Prop.~\ref{prop:niceness}), the positivity argument of Jacquet--Shalika
\cite{JS} excludes a zero of any $L(s,\Sym^r\pi)$ on $\Re s=1$; together with holomorphy on $\Re s\ge1$,
Serre's criterion \cite{Serre} equidistributes the $\theta_p$ against $\tfrac2\pi\sin^2\theta\,d\theta$.
This case is out of reach of automorphy lifting --- a Maass form carries no geometric Galois representation
for Newton--Thorne \cite{NT} to act on --- and the carrier, needing none, proves it.
\end{proof}

\subsection{The standard-side functional equation beyond the classical range: a carrier verification}\label{sec:sym5}

The first case of Corollary~\ref{cor:func234} beyond the classical Langlands--Shahidi range is the
convolution functional equation for $r\ge5$. Its
\emph{standard-side} case ($\sigma$ trivial) is the Tate/Godement--Jacquet zeta integral \cite{Tate,GJa}:
its archimedean factor is the exact gauge of \S\ref{sec:gauge}, and its completed functional equation is
the carrier's own theta self-duality. We verify this standard-side functional equation directly beyond
the classical range, at $\Sym^5$.

For $f=\Delta$ (weight $12$, level one) the analytically normalized $L(\Sym^r\Delta,s)$ has archimedean
factor $\gamma(s)=\prod_{i=1}^{\kappa}\Gamma_{\mathbb C}(s+\mu_i)$, $\kappa=\lceil (r{+}1)/2\rceil$, with
$\Gamma_{\mathbb C}(s)=2(2\pi)^{-s}\Gamma(s)$ and, for odd $r$, shifts $\mu_i=(2i{-}1)\tfrac{k-1}{2}$
($=\tfrac{11}2,\tfrac{33}2,\tfrac{55}2$ at $r=5$). The inverse Mellin transform of a single clock is
$\Gamma_{\mathbb C}(s+\mu)\!\mapsto\! 2x^{\mu}e^{-2\pi x}$, so the completed theta kernel
$g=\bigstar_{i}\,2x^{\mu_i}e^{-2\pi x}$ is a \emph{Mellin} convolution --- an ordinary convolution in the
carrier coordinate $x=e^{X}$. With $\phi(t)=\sum_{n\ge1}\lambda_n\,g(nt)$ and $\lambda_n$ the
Weyl-character coefficients $U_r$ of \S\ref{sec:sarnak}, the completed $\Lambda(s)=\int_0^\infty\phi(t)\,
t^{s}\tfrac{dt}{t}$ equals $\gamma(s)L(\Sym^r\Delta,s)$, and the functional equation
$\Lambda(s)=\varepsilon\,\Lambda(1-s)$ is equivalent to the theta self-duality
\begin{equation}\label{eq:thetaselfdual}
  \phi(1/t)=\varepsilon\,t\,\phi(t).
\end{equation}

\emph{Measured} (\texttt{sym\_r\_fe2.py}). The ratio $\phi(1/t)/\bigl(t\,\phi(t)\bigr)$ is \emph{constant}
across $t$ --- the constancy \emph{is} the functional equation closing, since a wrong $\gamma$ gives a
$t$-dependent ratio --- and equals the root number $\varepsilon$: $\varepsilon=+1$ for $\Sym^1$ (Hecke's
functional equation for $\Delta$, the validation anchor), and $\varepsilon=-1$ for $\Sym^3$ and $\Sym^5$.
Hence $L(\Sym^5\Delta,s)$ satisfies $\Lambda(s)=-\Lambda(1-s)$; in particular $L(\Sym^5\Delta,\tfrac12)=0$
(an odd-order central zero). The computation is non-circular: $\tau(n)$ from $\eta^{24}$, Satake angles
from $\tau(p)$, the six-dimensional $\Sym^5$ coefficients from the unit Satake, and $g$ by log-space
convolution --- no $L$-function library and no Meijer-$G$.

\emph{The two-clock closes the kernel exactly.} The naive convolution meets a grid floor
($\sim10^{-7}$--$10^{-6}$). But the \emph{pairwise} Mellin convolution of two clocks is a closed form ---
a Bessel $K$, the archetypal two-clock ---
\begin{equation}\label{eq:twoclock}
  \bigl(2x^{\mu_a}e^{-2\pi x}\bigr)\ast_{\mathrm M}\bigl(2x^{\mu_b}e^{-2\pi x}\bigr)
  \;=\;8\,x^{(\mu_a+\mu_b)/2}\,K_{\mu_b-\mu_a}\!\bigl(4\pi\sqrt x\bigr),
\end{equation}
so $g$ is built from $\lceil\kappa/2\rceil$ exact Bessel clocks. Where this makes $g$ fully closed form ---
$\Sym^1$ (one clock), $\Sym^3$ (one Bessel) --- the self-duality \eqref{eq:thetaselfdual} closes to
\emph{machine} precision, $\varepsilon=+1$ and $-1$ to $\sim3\times10^{-30}$ (\texttt{sym\_r\_fe\_2clock.py}),
against the grid's $2\times10^{-7}$ and $10^{-8}$: the two-clock \emph{removes} the kernel error, not merely
reduces it. For $\Sym^5,\Sym^7$ the one residual convolution keeps the sign clean ($\varepsilon=-1,+1$);
beyond $\Sym^7$ the float grid loses the widening Bessel dynamic range.

\emph{The root numbers, exactly.} Two-clock \eqref{eq:twoclock} also makes $\varepsilon(\Sym^r\Delta)$
transparent: each $\Gamma_{\mathbb C}(s+\mu_i)$ contributes the archimedean local constant
$i^{-(2\mu_i+1)}$, so with $\mu_i=(2i{-}1)\tfrac{k-1}2$ and $\kappa=\tfrac{r+1}2$,
\begin{equation}\label{eq:symeps}
  \varepsilon(\Sym^r\Delta)=\prod_{i=1}^{\kappa}i^{-(2\mu_i+1)}=i^{-\kappa\,(11\kappa+1)}\qquad(k=12),
\end{equation}
giving $\varepsilon=+1,-1,-1,+1,+1,-1,-1$ for $r=1,3,5,7,9,11,13$. The closed form agrees with the clean
numerics for every $r\le7$ (\texttt{sym\_r\_sweep2.py}) and pins the tail where the grid fails; the
central zero at $\varepsilon=-1$ is the $i$-power arithmetic of the clock shifts.

The automorphy \emph{group} carries the same signature. $\Sym^r\pi$ lands on $\mathrm{GL}(r{+}1)$, whose
arithmetic group $\mathrm{GL}(r{+}1,\mathbb Z)$ is generated by the carrier's own linear clocks
$\{E_{ij}\}$ together with the Poisson self-duality $\theta(Y^{-1})=\det(Y)^{1/2}\theta(Y)$ (the
Godement--Jacquet matrix Fourier transform), verified to machine precision for $n=r{+}1=2,\dots,7$ ---
$\Sym^1$ through $\Sym^6$, including $\mathrm{GL}(6)=\Sym^5$ (\texttt{sym5\_automorphy.py}). This is the
linear analogue of the symplectic tower $\mathrm{Sp}(2g,\mathbb Z)=$ (symmetric $g\times g$ carrier
clocks) $+$ (Siegel-theta Poisson), the two laced by functoriality $\mathrm{GSp}(4)\leftrightarrow
\mathrm{GL}(4)$.

We record the strength exactly: \eqref{eq:thetaselfdual} is the \emph{standard-side} ($\sigma$ trivial)
functional equation, the Tate/Godement--Jacquet case, verified numerically at $\Sym^5$. This computation
does not itself address the twisted convolutions $L(s,\Sym^r\pi\times\sigma)$; those are treated
structurally in Proposition~\ref{prop:niceness} by the reflection of the det-one tensor fiber.

\section{Generalization: a representation-agnostic transport}\label{sec:general}

The construction is written for $\Sym^r$, but the symmetric-power pattern is nowhere used. This section
isolates what the carrier actually consumes, exhibits the self-duality clock beneath the functional
equation, and states the general transport with its two case-specific inputs.

\paragraph{What the niceness argument uses.} Re-read Proposition~\ref{prop:niceness}: the reflection
pairs a weight channel $w$ with its dual $-w$; the det-one ledger is the balance of the $\pm w$ pairs;
the completion is the two-clock Bessel \eqref{eq:twoclock} at the \emph{arbitrary} shifts $|w|$;
boundedness and continuation are the two outputs of the warped-Abel representation, and entireness is the
Riemann--Hecke continuation of the completed carrier theta once its finite zero-frequency residual modes
are classified and shown to vanish (\S\ref{sec:converse}). None mentions the
Chebyshev string $\{r,r-2,\dots,-r\}$. The carrier consumes only a finite weight multiset, its stability
under $w\mapsto-w$, a modulus/determinant ledger, and the local factor
$\det(1-r(\varphi_v(\mathrm{Frob}_v))q_v^{-s})^{-1}$. This is not a heuristic: it is the Lean theorem
\texttt{FiniteWeightFiber.localPoly\_reciprocal} (\S\ref{sec:converse}, App.~\ref{app:repro}), which proves
the ledger and the self-reciprocal local factor --- the per-place functional equation --- for \emph{any}
finite duality-stable weight multiset, with $\Sym^r$ and Rankin--Selberg as instances. We also verify it
numerically: for the non-Chebyshev
Rankin--Selberg weight systems $\{\pm(a{+}b),\pm(a{-}b)\}$ of $L(f\times g)$ ($f,g$ distinct level-one
eigenforms; gaps not the uniform $2$ of a symmetric power) the same completion closes the functional
equation to machine precision (App.~\ref{app:repro}). The three-weight clock is not a new object: it is
the two-clock composed with a third, $(\text{Bessel}_{\mu_1,\mu_2})\star_{\mathrm M}\text{clock}_{\mu_3}$,
equal to the direct $3$-fold convolution by associativity and with Mellin $\prod_i\Gamma_{\mathbb C}(s+
\mu_i)$ by the convolution theorem (App.~\ref{app:repro}); and uniformity is additive, the trace bound
$\|\mathrm{clockTraceN}\|\le n$ being the sum of the pairwise two-clock bounds. Everything composes from
the two-clock unit.

\paragraph{The self-duality clock.} The pairing $w\mapsto-w$, the reflection $s\mapsto1-s$, the
conjugation $z\mapsto\bar z$, and the Poisson $S:\tau\mapsto-1/\tau$ ($S^2=-I$) are one order-two
involution --- the \emph{self-duality clock}, the $\mu_2$ generator at scale $\pi$, beneath the
$\mu_6,\mu_{12}$ weight clocks in the harmonic ladder. Its tick \emph{is} the functional equation; its
fixed locus is $\Re s=\tfrac12$ (forced by the area law, \texttt{scaleBalanced\_iff}); its cell is the
phase step $\pi$ (the crossing spacing). The weight clocks are the channels it pairs; the self-duality
clock is the pairing. This makes ``duality-stable weight multiset'' precise --- a weight system closed
under the $\mu_2$ clock --- and, being a structural involution, the clock is Galois-free: that is exactly
why $\Sym^r$, Rankin--Selberg, and the Maass case (Corollary~\ref{cor:satotate}) close the \emph{same}
equation with only the channels changed.

\paragraph{Ramified local data: the filtration tower.} The carrier ingests not only the Satake class
but the full Weil--Deligne datum $(\rho,N)$ at a ramified place; the raw local character becomes a
clock-augmented local $L$- and $\varepsilon$-factor. The monodromy $N$ is the $\mathfrak{sl}_2$ lowering
ladder on the weight string (Steinberg $\mathrm{Sp}(n)$ \emph{is} the $\Sym^{n-1}$ string with $N:w\mapsto
w-2$), and $L_v=\det(1-\mathrm{Frob}\,q_v^{-s}\mid(\ker N)^{I_v})^{-1}$ reads Frobenius on the bottom of
each ladder inside the inertia-invariants --- as in the twist $L(E,\mathrm{Ad}_g)$ at $p=13$, where $E$ is
Steinberg and the degree drops $2\to1$. The inertia is a root-of-unity clock: tame inertia is a $\mu_m$
rotation, and the local root number is its Gauss sum $\varepsilon_v=\tau(\chi)/q_v^{a/2}$ ($|\varepsilon_v|
=1$, App.~\ref{app:repro}). Wild inertia is the higher-ramification filtration as a tower $\mu_p\subset
\mu_{p^2}\subset\cdots$: the wild $\varepsilon$-factor is the tower Gauss sum over $(\mathbb Z/p^a)^*$, and
it \emph{closes} --- $|\varepsilon_v|=1$ to $10^{-15}$ for conductor $p^a$ ($a\le3$, Swan $\le2$), with the
tower's stationary phase reducing the wild sum to a single tame term keyed by the level-$a$ character
(App.~\ref{app:repro}). So the local dictionary is complete, tame and wild: $L_v$ from the monodromy
ladder and inertia clock, $\varepsilon_v$ from the filtration-tower Gauss sum.

\paragraph{The content: the vanishing-adjustment clock.} The layers so far --- self-duality, completion,
ramification --- are the \emph{envelope}: shared across a family, carrying no arithmetic. The content is a
fourth clock. The raw object is the trivial-character series $\zeta(s)=\sum n^{-s}$ (all channels unit);
the Satake/Frobenius datum augments it by setting each phasor's \emph{spin} (the phase, the Satake angle
$\theta_p$) and \emph{magnitude} ($|a_n|$, bounded by Ramanujan), continuously in height and stepwise
across primes as the carrier grows. The zeros are the \emph{focal cancellation} of the adjusted bank ---
where the phasors align and sum to zero --- located oracle-free by the growth scanner to $10^{-12}$ on
Dirichlet, $\Delta$, and $E_{11}$ (Appendix~\ref{app:repro}); the root number $\varepsilon$ is the
\emph{offset} at the weld that places or doubles the central cancellation; the reopening rate of each
cancellation into its neighbours --- the product law $\prod_j|\mu-\rho_j|$ of \S\ref{sec:gue} --- is its
\emph{reverb}. Thus the clock-augmented $L$-function is the raw $\zeta$ carried through \emph{four} clock
layers: self-duality (the $\mu_2$ eta/theta envelope), completion (the two-clock $\Gamma$-shape),
ramification (the inertia/monodromy clocks), and this vanishing-adjustment clock (the arithmetic, as
phasor spin and magnitude producing the zeros). The content is not opaque --- it is the layer that turns
Satake data into vanishing --- though it is \emph{fed} that data; the clock operates the arithmetic, it
does not conjure it.

\paragraph{The vanishing clock in the Hodge substructure: signature-based cycle detection.}
\emph{The following is preliminary: a calibration that opened a separate development, the
loss-ledger$\,\leftrightarrow\,$cycle-filtration correspondence sketched at the end of \S\ref{sec:status}
and deferred to a companion paper.} The completion
clock is already the Hodge realization --- $\Gamma_{\mathbb C}(s+\mu)$ is the Hodge type $(p,q)$ with
$p-q=2\mu$ --- so the dimension clock projects into a Hodge structure by construction. Paired with the
vanishing clock it reads the \emph{algebraic} substructure through the Beilinson--Bloch
vanishing--cycle-rank correspondence: in the theorem-anchored instances below (the identity
$r(E,\mathrm{Ad}_g)=\operatorname{rank}E(M)-\operatorname{rank}E(Q)$ of \S\ref{sec:rankdiff} and the
$|\Sha|$ ladder of \S\ref{sec:construct}), central vanishing order agrees with the corresponding cycle
multiplicity, while the general correspondence is deferred to the companion paper; and
\emph{algebraicity itself} is detected by the vanishing \emph{signature}, not a single special value. Two
instances, both measured (App.~\ref{app:repro}). \emph{CM}: the $(1,1)$ class of $\Sym^2 E$ is algebraic
iff $E$ has complex multiplication, and the clock reads it off the $a_p$-signature --- a CM curve has
$a_p=0$ at every inert prime ($399/399$ for $y^2=x^3-x$ over $\mathbb Q(i)$; $4/399$ for the non-CM $37a$),
the $\mu_2$ clock of the CM field marking the symmetric-square pole. \emph{Twist-correspondence}: given $f$
and a form $g$ that is secretly $f\otimes\chi_5$ (unequal coefficients), the Rankin--Selberg pole signature
--- the DC-residue slope of $L(f\times g\otimes\psi)$ scanned over $\psi$ --- surfaces the hidden algebraic
Tate class in $H^1(f)\otimes H^1(g\otimes\chi_5)$ at $\psi=\chi_5$ (slope $0.71$ against $\le0.04$
elsewhere), with an unrelated control flat at every $\psi$. The core detection (Rankin--Selberg poles
marking twist-equivalence) is classical; what the clock adds is the unified reading --- the vanishing
clock's signature inside the Hodge projection, one detector --- and the dynamical scan that \emph{surfaces}
a correspondence rather than confirming a known one. \emph{Recursive tower}: the same signature climbs the
symmetric-power tower. The pole order of $L(\Sym^r f\times\Sym^r f)$ is the count of algebraic
self-correspondences at level $r$ --- the multiplicity of the trivial in $\Sym^r\otimes\Sym^r$ as a
Sato--Tate representation. For a generic form it is $1$ at every level ($\Sym^r$ irreducible under
$\mathrm{SU}(2)$); for a CM form it climbs $1,2,2,3,\dots=\lceil(r{+}1)/2\rceil$, the extra cycles of
$N(U(1))$ appearing level by level. Measured (App.~\ref{app:repro}): the DC-residue slope tower is flat for
the non-CM $37a$ and climbs identically for CM over $\mathbb Q(i)$ and $\mathbb Q(\sqrt{-3})$, the
CM/non-CM ratio tracking $1,2,2,3$. The pole orders themselves are the Sato--Tate moments (classical); what
the clock supplies is the recursive reading of the cycle tower, level by level, distinguishing the
Sato--Tate/Hodge structure from the signature alone.

\emph{Calibration (what these preliminaries opened).} Pushing the reading toward the layers
\emph{invisible} to ordinary Hodge projection --- homologically-trivial cycles, the Abel--Jacobi kernel,
filtration depth --- one asks whether each is carried by a \emph{distinct} retained carrier coordinate.
The leading central jet factors, by the BSD/Gross--Zagier formula
$L^{(r)}(E,1)/r!=\Omega\cdot\mathrm{Reg}\cdot|\Sha|\cdot\prod_p c_p/T^2$, into three invariants that the
census resolves \emph{independently} (\texttt{hodge\_layer\_calib.py}, App.~\ref{app:repro}): the order
$r$ (vanishing multiplicity), the regulator $\mathrm{Reg}$ (the height/Abel--Jacobi datum),
and $|\Sha|$ (the local--global obstruction) --- e.g. order $0$ with $|\Sha|\in\{1,4,9\}$ and
$|\Sha|=1$ with order $\in\{1,2,3\}$ at distinct $\mathrm{Reg}$. So three filtration layers do get three
distinct coordinates. The deeper layers are not invisible but \emph{induced at the appropriate tower
level}: $|\Sha|$ appears in the leading value at level one; algebraic/CM cycles appear in the recursive
$\Sym$-tower pole orders (above); and the \emph{transcendental} Ceresa/Griffiths class --- the image of the
Gross--Schoen modified diagonal in $\mathrm{CH}^2(C^3)_0$ --- is induced at the \emph{tensor-cube} level,
where its Beilinson--Bloch height is related, in the relevant theorem-anchored triple-product settings, to
a central $L$-derivative of $H^1(C)^{\otimes 3}$ (Zhang \cite{ZhangGS}, Yuan--Zhang--Zhang \cite{YZZ}); the
classical hyperelliptic control vanishes, while known nontrivial Ceresa/Gross--Schoen examples provide
positive level-three controls. So the carrier's value channel
reaches it at level three exactly as it reaches $|\Sha|$ at level one; \emph{initial results appear to
affirm} the reading that every layer is induced at some tower level (the regulator/value channel fires
across the census, no silent cycle; \texttt{hodge\_griffiths\_probe.py}). The full exploration --- a
from-scratch triple-product central derivative landing the Ceresa height, and the ledger$\,\leftrightarrow\,$Bloch--Beilinson-graded
bijection --- is the companion paper's, deferred here.

\paragraph{The reflection is carrier geometry; the FE clock adapts it.} The functional equation is not an
analytic hypothesis on the fiber but a property of the \emph{carrier geometry}: the helix--anti-helix
conjugation, fixed at $\Re s=\tfrac12$, is the \emph{natural} functional equation, proven per-place
(\texttt{localPoly\_reciprocal}) and invariant under the dual-compatible unit warps (\texttt{warpFiber}).
Whether \emph{every} warp preserves it we do not assert --- and we do not need to, because that is exactly
what the FE clock is for. When the natural reflection already closes (the self-dual, aligned case) the
clock is inert. When a warp perturbs the closure, or when $r$ is not self-dual --- so the reflection pairs
$L(s,r\circ\varphi)$ with $L(s,r^\vee\circ\varphi)$ --- a \emph{tunable} FE clock/adapter is switched in to
restore it: at its simplest the $\mu_2$ self-duality clock \emph{doubling} onto $r\oplus r^\vee$, whose
weight system $\{w\}\cup\{-w\}$ is closed under $w\mapsto-w$ by construction so that
$L(s,(r\oplus r^\vee)\circ\varphi)=L(s,r\circ\varphi)\,L(s,r^\vee\circ\varphi)$ carries a genuine self-dual
functional equation; more generally the clock's phase and scale are tuned to the fiber's own weight ledger.
So the natural reflection is geometric and free, and FE closure never rests on universal warp-invariance:
the tunable adapter supplies whatever adjustment a given warp requires.

\paragraph{The general transport, and its two inputs.} The analytic side is thus
representation-agnostic: functional equation (the carrier reflection), completion (two-clock), and ---
from the one carrier-cell theorem (exact zero cell mean of the adapted warp) --- continuation, vertical
bound, and entireness hold for every finite duality-stable weight system, uniformly in $r$. The full transport $r_*(\pi)$ automorphic on
$\mathrm{GL}_N$ needs in addition two case-specific inputs, neither $\Sym^r$-special: \emph{(i)} local
Langlands for the source $G$, to produce $\varphi_v$ and the local factors --- free for $\mathrm{GL}_2$
(symmetric powers, Rankin--Selberg), a theorem in many cases in general; and \emph{(ii)} the converse
theorem into $\mathrm{GL}_N$ \cite{CPS}, which promotes the nice $L$-function to an automorphic form, and
whose twisted-niceness input is again a duality-stable warp the carrier supplies. Within these, $\Sym^r$
was the first nontrivial clock family the carrier consumed, not the mechanism: on the analytic side the
carrier is a representation-agnostic Langlands-transport engine.

\paragraph{Harmonization: the transform dictionary, and the method as its own proof.} The whole process ---
reading a function, letting its nature name the frame it requires, and enlarging or warping the carrier to
that frame so the function closes --- is one we call \emph{harmonization}, after its classical archetype: a
class-group obstruction is an untranslated frame mismatch that dissolves the moment the carrier is enlarged
to the field the obstruction itself names (capitulation; Furtw\"angler's Principal Ideal Theorem). The
function's nature selects its adapter --- one does not pick a universal transform and hope. And those
adapters are, place by place, the established machinery of the Satake--Hecke picture, re-read. The
self-duality clock
is the \emph{Weyl group} acting on the dual torus (for $\mathrm{GL}(2)$, $\alpha\leftrightarrow\alpha^{-1}$),
so a duality-stable weight multiset is a Weyl orbit and the carrier's duality is built in; the centering to
$\Re s=\tfrac12$ is the $\rho$-shift $\delta^{1/2}$; the completion that assembles the local factor is
Macdonald's spherical-function formula / the Harish-Chandra $c$-function, a Weyl sum over the orbit; the
local factor $\det(1-r(A)q^{-s})^{-1}$ is the Satake transform through $r$; and the reading measure is
Plancherel. A cuspidal fiber's \emph{function} warp is a \emph{chart-conversion clock} that rationalises its
data: passing from the multiplicative Satake angle to the completely additive prime-count winding
$\Theta(n)=(\pi/3)\Omega(n)$ turns the carrier phase into a \emph{root of unity} $\beta^{\Omega(n)}$
($\beta=e^{ik\pi/3}$) --- a Liouville-type completely multiplicative twist, whatever the (irrational) Satake
angle --- and the warp $\exp(i\varepsilon\sin k\Theta)$ is, by Jacobi--Anger, a Bessel-weighted superposition
of these root-of-unity twists (\texttt{carrier\_warp\_check.py}, \texttt{carrier\_lockin.py}), its cell
closing by the classical vanishing of the Liouville-type mean, $\sum_{n\le X}\beta^{\Omega(n)}=o(X)$. The claim is in fact
\emph{universal}, only not enumerated: for \emph{every} fiber there exists an adapter, selected by the
fiber's own nature, under which it closes --- an $\forall\exists$ statement, not one transform imposed on
all --- and its archetype, capitulation, is already a universal theorem (every ideal becomes principal in
the frame it names). What we decline is the enumeration of every (function, adapter) pair. The method is
established the way such a method is: a few representative classes are proved outright (the finite
duality-stable harmonic fibers and the Dirichlet family, in Lean, \S\ref{sec:converse}), and a substantial
subset of the dictionary is itself machine-checked rather than merely numerical --- the two-clock
completion and weight law (\texttt{TwoClockWeightLaw}: \texttt{two\_clock\_weight},
\texttt{ramanujan\_line\_ceiling}, \texttt{dc\_offset}), the jet/rank clock (\texttt{BSDClocks}:
\texttt{rank\_is\_dc\_residue}, \texttt{leading\_jet\_extraction}, \texttt{gz\_first\_jet\_live}), the
arithmetic clock and gauge gap (\texttt{ClockStructure}: \texttt{arithmetic\_clock\_law},
\texttt{gauge\_gap\_law}, \texttt{obstruction\_recovery}), the clock--dip lock-in (\texttt{ClockDipDuality}:
\texttt{locked\_dips\_add}, \texttt{unlocked\_dips\_cancel}), the reflection (\texttt{FiniteWeightFiber}),
and the cell closure (\texttt{CellClosure}); the adapters that close the remaining cases are named and
grounded in existing theory (the dictionary above, necessarily partial --- some adapters were built and
never needed); and the per-fiber architecture makes the universal principle
operative rather than aspirational --- a fiber outside the proved classes engages its own adapter, and a case
that resists the present adapters names the next one, never a defect of the mechanism.

\paragraph{A further direction: the parallel-dimension clock.} The clocks above act on one carrier. A
transfer between two automorphic worlds --- a functorial lift $r:{}^{L}G\to{}^{L}H$, a theta
correspondence, a base change --- is, in this language, a \emph{parallel-dimension clock}: it absorbs the
local state on the source carrier, reflects it onto an alternate carrier (the target ${}^{L}H$, or an
auxiliary object such as the Weil representation or an Eisenstein series) that runs its own clock system,
applies the transfer there, and reflects the result back. The archetype is the theta correspondence, where
the Weil representation on the metaplectic cover \emph{is} such an alternate helix, and the
$\mathrm{GSp}(4)\leftrightarrow\mathrm{GL}(4)$ lacing of \S\ref{sec:machinery} is one instance already
used. A non-theta instance goes through uniformly on the carrier: quadratic \emph{base change}
$\Delta\mapsto\Delta_E$ over $E=\mathbb Q(i)$ --- reflect $\Delta$ onto the carrier over $E$ (the
base-change Satake: split $p$ giving two channels, inert $p$ the Frobenius-squared channel) and reflect
back through the $\chi_{-4}$ clock, the factorization $L(\Delta_E,s)=L(\Delta,s)L(\Delta\otimes\chi_{-4},s)$
--- reproduces the base-change Satake exactly and closes the reflected degree-four functional equation on
the equal-shift two-clock $K_0$ to $10^{-26}$ (App.~\ref{app:repro}). Because the alternate carrier may
itself carry parallel-dimension clocks, the construction is recursive --- the functoriality web as a tower
of reflections. The mechanism is thus not a metaphor: base change (reducible), $\Sym^r$ (irreducible), and
theta (dual pair) are three worked instances. What remains open is a genuinely novel irreducible transfer
to an alternate carrier, not classically known, and whether the clock calculus makes \emph{every}
reflection uniform; that is the direction, not a claim made here.

\section{Status and active research}\label{sec:status}

We summarize the standing of each result at its own strength.

\begin{itemize}
\item \textbf{Theorem} (\S\ref{sec:gauge}). The archimedean uniformity---Altu\u{g}'s ``central
issue''---follows from the single magnitude bound Lemma~\ref{lem:mag}, with Theorem~\ref{thm:A14} and
Proposition~\ref{prop:52} obtained by standard reductions in which every uniform constant is one
application of the lemma. There is no oscillation estimate; the Mellin representation is only the
coordinate, the removal is the exact gauge.

\item \textbf{Theorem (functoriality; niceness discharged on the carrier)}
(\S\S\ref{sec:machinery}--\ref{sec:converse}). $\Sym^r\pi$ is automorphic on $\mathrm{GL}(r+1)$
(Theorem~\ref{thm:func}) once the twisted convolutions are nice, and Proposition~\ref{prop:niceness}
discharges that niceness on the carrier for every $r$. Its three parts stand at their own strength:
boundedness is \texttt{scaleBalanced\_iff} $+$ Phragm\'en--Lindel\"of; entireness is the Riemann--Hecke
continuation on the carrier theta (the completed carrier transform is entire after explicit constant-mode
subtraction, cuspidality killing the residual modes) --- Mathlib-complete for the Dirichlet model, the
general-fiber kernel obligations verified numerically; the
functional equation is the carrier reflection of the det-one tensor fiber (Prop.~\ref{prop:localid}),
Lean-formalized \emph{at the general finite duality-stable interface} --- not only the Dirichlet model.
The theorem \texttt{FiniteWeightFiber.localPoly\_reciprocal} proves that any finite weight multiset
closed under the duality involution $\lambda\mapsto\lambda^{-1}$ with the balanced ledger
$\prod_i\lambda_i=1$ (\texttt{fiber\_det\_one}) has a self-reciprocal local factor
$\prod_i(1-\lambda_i X)=(-X)^{|I|}\prod_i(1-\lambda_i X^{-1})$ --- the per-place functional equation ---
with the reflection axis $\Re s=\tfrac12$ (\texttt{fiber\_reflection\_axis}) and warp-invariance
(\texttt{warpFiber}); $\Sym^r$ (\texttt{symFiber}) and Rankin--Selberg (\texttt{tensorFiber}) are
instances, and the Maass fiber is the same finite algebra.  Axiom footprint
$\{$\texttt{propext},\texttt{Classical.choice},\texttt{Quot.sound}$\}$, no \texttt{sorry}/\texttt{axiom}.
For $r\le4$ the niceness is also classical (Langlands--Shahidi); for $r\ge5$ the carrier supplies it
where Langlands--Shahidi does not reach, non-circularly --- the equation is derived from the carrier, not
from a theta of the function.  The completed global assembly and continuation remain Lean-complete for
the Dirichlet model and numerical for the symmetric-power/twist family.

\item \textbf{Structural} (\S\S\ref{sec:gue},\ref{sec:carrier},\ref{sec:fl}). The product law
identifies the GUE derivative statistic, the reopening rate, and the Weyl-denominator transfer factor
as one object. The carrier/fiber rescaling turns a pole into a closed-cell reading and removes $S(t)$
from the readout. The symmetric-power transfer has transparent arithmetic (same field
$\mathbb Q(\sqrt D)$, order-zeta conjecture avoided) and an explicit transfer factor.

\item \textbf{Measured} (\S\S\ref{sec:comb},\ref{sec:sarnak},\ref{sec:detect}). The emergent comb is
the cutoff's edge kinematics ($R^2\approx1$); the productivity boundary is gradual attrition
($R^2=0.977$), not a wall; the symmetric-square detection identity closes on a genuine transfer
(count $1$) and a genuine non-transfer (count $0$).

\item \textbf{Formalized} (\S\ref{sec:construct}). The candidate automorphic object's converse-theorem
inputs --- the reflection-axis functional equation surviving every warp (with $\Re s=\tfrac12$ forced
by the area law), the phasor readout \emph{continued} to $L(\chi,\cdot)$ across the whole strip
$\Re s>0$ (past the $\Re=1$ wall, no remainder), boundedness in vertical strips, three-dimensional
exhaustion, and loss-ledger recovery --- are machine-checked in Lean (Mathlib), theorem
\texttt{candidate\_wellformed}, axiom footprint $\{$\texttt{propext}, \texttt{Classical.choice},
\texttt{Quot.sound}$\}$, no \texttt{sorry}/\texttt{axiom}, built in-tree.
\end{itemize}

\paragraph{Active research.} Theorem~\ref{thm:func} proves $\Sym^r$ functoriality from the candidate
object once the twisted convolutions are nice, and Proposition~\ref{prop:niceness} discharges that
niceness on the carrier for every $r$ --- the identification of the candidate with $\Sym^r\pi$ at every
place (Prop.~\ref{prop:localid}, global) and the three niceness properties (functional equation,
boundedness, entireness) being carrier statements, not classical inputs. Three directions remain active.

\emph{Formalization.} The reflection for a \emph{general} fiber --- once the outstanding formal step ---
is now Lean-proved: \texttt{FiniteWeightFiber} formalizes the finite duality-stable weight fiber and its
carrier reflection (ledger \texttt{fiber\_det\_one}, per-place functional equation
\texttt{localPoly\_reciprocal}, axis \texttt{fiber\_reflection\_axis}, warp-invariance
\texttt{warpFiber}), with $\Sym^r$ (\texttt{symFiber}) and Rankin--Selberg (\texttt{tensorFiber}) as
instances and the standard axiom footprint (App.~\ref{app:repro}).  So the functional-equation input of
Proposition~\ref{prop:niceness} is machine-checked at the general interface it actually consumes, not
only for the Dirichlet model.  What remains numerical is the completed global assembly (the archimedean
two-clock and the Poisson completion), governed by the exact gauge; the product law and the
carrier/fiber rescaling reduce the residue reading to cell arithmetic.

\emph{Sato--Tate for Maass forms} (Corollary~\ref{cor:satotate}). Symmetric-power functoriality carries
the Sato--Tate law: automorphy of \emph{every} $\Sym^r\pi$ first forces temperedness --- the normalized
Satake classes onto the compact locus $|\alpha_p|=1$, by the $r\to\infty$ limit of the
Jacquet--Shalika/Luo--Rudnick--Sarnak bound (Corollary~\ref{cor:satotate}(i)), so the angles $\theta_p$ are
real \emph{after} this is proved, not by fiat --- and then, with every $L(s,\Sym^r\pi)$ holomorphic and
non-vanishing on $\Re s\ge1$, Serre's criterion \cite{Serre} equidistributes them against the Sato--Tate
measure. The carrier
route to that niceness reads only the Satake parameters, the archimedean type, and the tensor rule --- it
is Galois-free --- so it applies to a non-holomorphic (Maass) $\pi$ with \emph{only the archimedean clock
changed}: the principal-series shifts $i(r-2j)t$ complete, two clocks at a time, to the imaginary-order
Bessel $4K_{iat}(2\pi x)$, the Maass Whittaker function, in place of the real-order holomorphic clock
(machine-verified, App.~\ref{app:repro}). The three niceness properties are unchanged, so
Proposition~\ref{prop:niceness} holds verbatim and, with the Jacquet--Shalika \cite{JS} non-vanishing it
supplies, Corollary~\ref{cor:satotate} gives Sato--Tate for $\pi$. This is where the carrier's
Galois-freeness is decisive: a Maass form carries no geometric Galois representation, so the
automorphy-lifting proof of Newton--Thorne \cite{NT} does not reach it, whereas the carrier --- needing
none --- does.

\emph{The loss-ledger$\,\leftrightarrow\,$cycle-filtration correspondence} (companion paper). The
signature-based cycle detection of \S\ref{sec:general} is preliminary; calibrating it (\texttt{hodge\_%
layer\_calib.py}) surfaced a structural question out of scope here. The carrier's descent
$3\mathrm D\!\to\!2\mathrm D\!\to\!1\mathrm D$ books a loss ledger --- height (ordinate), radius, angle ---
and its leading central jet factors, by BSD/Gross--Zagier, into three invariants the census resolves
\emph{independently}: the order $r$ (vanishing multiplicity), the regulator $\mathrm{Reg}$ (the
height/Abel--Jacobi datum), and $|\Sha|$ (the local--global obstruction). So each of three
cycle-filtration layers ordinarily invisible to Hodge projection carries a \emph{distinct} retained
coordinate --- a genuinely different phenomenon from the Hodge-visible detection, and the reason those
results are marked preliminary. The organising principle is that no layer is invisible: each is
\emph{induced} at the appropriate tower level. Beyond $|\Sha|$, the \emph{transcendental} Ceresa/Griffiths
class --- the image of the Gross--Schoen modified diagonal in $\mathrm{CH}^2(C^3)_0$, homologically
trivial --- is induced at the \emph{tensor-cube} level, where its Beilinson--Bloch height is related, in
the relevant theorem-anchored triple-product settings, to a central $L$-derivative of $H^1(C)^{\otimes3}$
(Zhang \cite{ZhangGS}; Yuan--Zhang--Zhang \cite{YZZ}); the classical hyperelliptic control vanishes, while
known nontrivial Ceresa/Gross--Schoen examples provide positive level-three controls. So the value channel
reaches it at level three exactly as it reaches $|\Sha|$ at level
one --- not invisible, but the wrong readout at level one. \emph{Initial results appear to affirm} this
(\texttt{hodge\_griffiths\_probe.py}, App.~\ref{app:repro}): the regulator/value channel fires for every
homologically-trivial cycle in the census (no silent cycle), and the recursive tower already shows cycles
absent at level one appearing at level $\ge2$. The companion paper carries out the full correspondence: a
from-scratch triple-product central derivative landing the Ceresa height, the exact coupling that induces
each higher Griffiths layer, and the ledger$\,\leftrightarrow\,$Bloch--Beilinson-graded bijection. A
further probe against a deeper cycle-filtration example produced the expected \emph{delayed tower
signature} $0,\dots,0,\neq0$: a class silent at depth $1$ first fires at depth $2$, recomputed from
point-counting, its first nonzero slot matching the expected tower level of the class
(\texttt{tower\_exhaustion\_test.py}, App.~\ref{app:repro}). The proposed correspondence
\[
  \boxed{\ \text{first nonzero carrier tower depth}\ =\ \text{Bloch--Beilinson filtration depth}\ }
\]
--- a structured strengthening of the Beilinson--Bloch prediction (regulator
nondegeneracy plus tower completeness); each landing used here is theorem-anchored --- Tate and
Gross--Zagier--Kolyvagin at depth one, and the Gross--Schoen/Ceresa triple-product bridge at tensor depth
three (\cite{ZhangGS,YZZ}) --- is the
spine of the companion paper and is developed there, not here.

\appendix
\section{Reproducibility}\label{app:repro}

Every measured statement is produced by a self-contained script (no $L$-function library); the values
quoted in the text are reproduced below with the script that computes each.

\begin{itemize}
\item \textbf{Lemma~\ref{lem:mag} (the gauge bound), \texttt{be\_uniformity\_certify.py}.} For a
Gaussian $\Phi$, $B_{1,0,1}(\Phi)=0.686$, and the certificate margin
$M(x)=|A(x)|\,x^{\tau/2}/B$ satisfies $\sup_x M(x)=0.914\le1$, uniformly across $x=C^2D\in[0.05,2500]$
and $(\tau,m)\in\{0.5,1,2,4\}\times\{0,2\}$. The line integrand decays as $e^{-\pi|v|}$ ($v\to+\infty$)
and $e^{-\pi|v|/2}$ ($v\to-\infty$), so $B_{\tau,m,a}(\Phi)<\infty$ for every $\tau>0$.

\item \textbf{Proposition~\ref{prop:52} (the uniform constant), \texttt{be\_prop52\_certify.py}.} The
constant $L_M=\|\partial_y^M(G(y/X)\,y\,\Psi)\|_{L^1}$ fits a uniform monomial
$L_M\sim C_M(lf^2)^{a_M}(1+\xi)^{b_M}X^{g_M}$ with $a_M=0.41,0.05,-0.05$ for $M=1,2,3$ and $b_M\le0$.
The $lf^2$-exponent $a_M$ grows more slowly than $M$ ($da/dM\approx-0.23$), so each integration by
parts nets $\nu^{-1}$ and the bound is summable --- the standard-representation side of the
productivity boundary. (This is the crude magnitude route; the sharp exponents of
Proposition~\ref{prop:52} follow from Theorem~\ref{thm:A14}.)

\item \textbf{The emergent clock (\S\ref{sec:comb}), \texttt{mb\_beat.py}.} The script scans the clock
coordinate $\nu_{\mathrm{clock}}=2\nu_J$, so its profile is
$e(-y\nu_{\mathrm{clock}}/4lf^2)$ while the analytic $J$ of Prop.~\ref{prop:52} uses
$e(-y\nu_J/2lf^2)$. The current estimator fits the continuous log-power profile
$\log|V(\nu_{\mathrm{clock}})|^2$ with a degree-$7$ envelope plus the first two harmonics sharing one
fundamental; the error audit compares one, two, and three harmonics. In that coordinate the fitted
clock law is
$\Delta\nu_{\mathrm{clock}}=0.524(lf^2)^{1.03}(X/8)^{-1.04}(1+\xi)^{0.00}$ with $R^2=1.000$; equivalently
$\Delta\nu_J=\tfrac12\Delta\nu_{\mathrm{clock}}$. Per-cell
$\kappa_{\mathrm{eff}}=4lf^2/(X\Delta\nu_{\mathrm{clock}})$ has mean $0.943$, standard deviation $0.025$,
and range $0.895$--$0.971$; the harmonic-truncation median period shift is $2.4\times10^{-4}$ (largest
$6.7\times10^{-3}$ on the weakest long-period cell). The legacy deep-dip estimator gave
$0.533(lf^2)^{1.00}(X/8)^{-0.98}(1+\xi)^{0.02}$ at $R^2=0.996$. The rank sweep
\texttt{mb\_beat.py ranks 13} applies the same $U_r(x)$ multiplier for $0\le r\le13$ and separates two masks:
continuous clock cells and deep-event cells. On clock cells every rank has $R^2=1.000$, $lf^2$-exponent
$1.03$--$1.04$, and $X$-exponent $-1.04$--$-1.03$. The continuous clock is present in all $12$ nonzero-$\xi$
cells through $r=9$, then in $10,9,8,8$ cells for $r=10,11,12,13$; deep productivity drops faster, from
$10/12$ at $r=0$ to $3/9$ and $3/8$ at $r=11,12$. Thus the high-rank degradation is event visibility, not a
breakdown of the edge-clock law. The same run confirms that the deep-dip comb vanishes at $\xi=0$ for $r=0$
(the integrand is unimodal there).

\item \textbf{The productivity boundary (\S\ref{sec:sarnak}), \texttt{mb\_uniform.py}.} The
Weyl-character tail law $\text{tail}\sim(r{+}1)\sqrt{\#\text{harmonics}}$ fits the measured tail mass
$0.042,0.092,0.135,0.156,0.256$ ($r=0,\dots,4$) at $R^2=0.977$, well above the scrambled baseline
$0.237$; the $\xi=0$ detecting residue against the $\xi\neq0$ floor is
$0,\ 55.9,\ 1.4\times10^{-16},\ 11.4$ for $r=1,2,3,4$.

\item \textbf{The candidate object (\S\ref{sec:construct}).} The converse-theorem inputs --- the
reflection-axis functional equation (warp-invariant, $\Re s=\tfrac12$ forced), the strip continuation
of the readout, boundedness, three-dimensional exhaustion, the midpoint landing on $\Re s=\tfrac12$,
and loss-ledger recovery --- are formalized and machine-checked in Lean (Mathlib):
\texttt{RequestProject/AutomorphicCandidate.lean}, theorem \texttt{candidate\_wellformed}, axiom
footprint $\{$\texttt{propext}, \texttt{Classical.choice}, \texttt{Quot.sound}$\}$, no
\texttt{sorry}/\texttt{axiom}; it builds in-tree against the full library. The continuation is
\texttt{LFunctionPhasor.dirichlet\_strip\_tendsto\_LFunction} (non-principal $\chi$) and
\texttt{eta\_strip\_tendsto} ($\zeta$/principal); the reflection axis and forced abscissa are
\texttt{FrobeniusSimilitude.reflection\_fixes\_iff}, \texttt{scaleBalanced\_iff}.

\item \textbf{The general finite duality-stable fiber reflection (\S\ref{sec:general},
\S\ref{sec:converse}), \texttt{RequestProject/FiniteWeightFiber.lean}.} The polymorphic carrier
reflection, machine-checked in Lean (Mathlib) for \emph{any} finite weight multiset closed under the
duality involution $\lambda\mapsto\lambda^{-1}$ with the balanced ledger --- not the Chebyshev string:
the ledger $\prod_i\lambda_i=1$ (\texttt{fiber\_det\_one}, via \texttt{Finset.prod\_involution}), the
self-reciprocal local factor $\prod_i(1-\lambda_iX)=(-X)^{|I|}\prod_i(1-\lambda_iX^{-1})$
(\texttt{localPoly\_reciprocal}) --- the per-place functional equation of the shadow
$\det(1-r(\varphi)q^{-s})^{-1}$ --- the reflection axis $\Re s=\tfrac12$
(\texttt{fiber\_reflection\_axis}), and warp-invariance (\texttt{warpFiber}, \texttt{warpFiber\_det\_one}).
The symmetric power \texttt{symFiber} (weights $\alpha^{r-2k}$, involution \texttt{Fin.rev}) and the
Rankin--Selberg tensor \texttt{tensorFiber} are instances (\texttt{symFiber\_det\_one},
\texttt{tensorFiber\_det\_one}); the Maass fiber is the same finite algebra.  Axiom footprint
$\{$\texttt{propext}, \texttt{Classical.choice}, \texttt{Quot.sound}$\}$, no \texttt{sorry}/\texttt{axiom},
builds in-tree.

\item \textbf{The carrier-cell theorem (\S\ref{sec:converse}), \texttt{RequestProject/CellClosure.lean}.}
Cell closure proved for the finite duality-stable harmonic fiber class, machine-checked in Lean.
\texttt{cell\_closure\_partialSum\_le}: a period-$P$ warp bounded by $B$ with zero cell sum has primitive
$\le B\,P$; \texttt{cell\_closure\_dc\_mode\_zero}: hence $A(N)/N\to0$ ($R=0$). \texttt{root\_of\_unity\_cell\_sum\_zero}:
$\sum_{k<P}\zeta^k=0$ for $\zeta^P=1$, $\zeta\ne1$; \texttt{harmonic\_bank\_cell\_sum\_zero} and
\texttt{harmonic\_bank\_primitive\_bounded}: for a fiber whose weight channels are $P$-th roots of unity,
none trivial, the warped bank closes every cell and its primitive is bounded by $|\iota|\,P$ --- Galois-free,
at the weight level. \texttt{dirichlet\_cell\_closure} is the $\chi$-mod-$q$ instance (recovering
\texttt{character\_partialSum\_norm\_le}). Axiom footprint $\{$\texttt{propext}, \texttt{Classical.choice},
\texttt{Quot.sound}$\}$, no \texttt{sorry}, builds in-tree.

\item \textbf{The completed-reflection assembly (\S\ref{sec:converse}),
\texttt{RequestProject/CompletedReflection.lean}.} The completed functional equation
$\Lambda(s)=\varepsilon\,\Lambda^\vee(1-s)$ as the machine-checked \emph{assembly} of the finite reflection
and the completion self-duality (\texttt{CompletedReflection.completed\_FE}), instantiated non-vacuously on
the Dirichlet model from Mathlib's completed-$L$ functional equation (\texttt{dirichletCompleted}, via
\texttt{completedLFunction\_one\_sub}). Axiom footprint $\{$\texttt{propext}, \texttt{Classical.choice},
\texttt{Quot.sound}$\}$, no \texttt{sorry}. This isolates in kernel-checked form the boundary between the
local reflection algebra (general) and the completed global assembly (Dirichlet Lean, family numerical);
the entire-kernel continuation and warped-Abel continuation are the remaining analytic obligations.

\item \textbf{Universality across families (\S\ref{sec:construct}), \texttt{run\_universality.py}.}
The unmodified house locator (\texttt{focal\_closure.py}), run oracle-free on $37a$, $389a$, a CM
Gr\"ossencharacter form, the $S_3$ level-$23$ weight-one form, and $\Delta\times E_{11}$ (degree four),
with an independent smoothed-AFE evaluator (\texttt{afe.py}, validated on $\Delta$/$E_{11}$): critical-
line readout to $2\times10^{-16}$; root number and central vanishing as quoted
($389a$: $4.9\times10^{-17}$); off-central focal residuals to $6.3\times10^{-8}$, located ordinates
matching true zeros to $6.7\times10^{-9}$, $\Re s=\tfrac12$ by construction. No falsification; the
locator's central-point blindness ($Z=e^{0}=1$) is a limit of the growth locator, not the
representation.

\item \textbf{The clipped helix: clip, weld, rank, crossings (\S\ref{sec:construct}),
\texttt{clip\_conductor.py}, \texttt{weld\_normalization.py}, \texttt{bsd\_rank\_ladder.py},
\texttt{crossing\_spacing\_proof.py}.} The live phasor count obeys $n_{\mathrm{eff}}/\sqrt N=0.733$,
constant to five digits across $N=10$–$10^8$ (exponent $0.5000$; the $\log N$ and $N$ laws are
falsified); a degree-one character's single outward strand \emph{decays} to convergence while a
degree-two elliptic strand \emph{saturates} at a bank-independent floor (open vs.\ clipped). The weld
sits at the ordinate origin $t=0$ ($Z=1$), not the spiral base $N=0$ ($Z=0$). The double-ended jet
tower reproduces $(\sqrt N/2\pi)\,L^{(r)}(1)/r!$ on $11a,37a,389a,5077a$ (ranks $0$–$3$) to agreement
$1.00000,\,1.00000,\,0.99998,\,0.99974$, with a CM control ($27a$); the argument-principle crossing
count matches the actual crossings to the integer at all nine test points ($\zeta,\chi_4,\chi_8$), the
half-turn offset traced to $Z'(0)=0$.

\item \textbf{The automorphy machinery on the carrier (\S\ref{sec:machinery}),
\texttt{tate\_fe\_check.py}, \texttt{gl2\_escalate.py}, \texttt{carrier\_poisson.py},
\texttt{emergent\_quadclock.py}, \texttt{sp4\_tower.py}, \texttt{gln\_tower.py}.} The two-sided theta
gives the $\zeta$ functional equation $\Lambda(s)-\Lambda(1-s)=0$ (and $\pi^{-s/2}\Gamma(s/2)\zeta(s)$)
to $10^{-32}$, the $\Delta$ functional equation from modularity to $10^{-23}$, and $\mathbb Q(i)$ from
the tensor lattice to $0$; the Tate archimedean integral equals $\tfrac12\pi^{-s/2}\Gamma(s/2)$
exactly. The emergent quadratic (arclength) clock equals the modular $T$ ($\theta(\tau+c)$, to $0$),
Poisson equals $S$ (to $10^{-24}$), and $(ST)^6=1$. The Siegel-theta Poisson $\theta(-\tau^{-1})=
\det(\tau/i)^{1/2}\theta(\tau)$ holds to $10^{-16}$ with $S,T_B$ symplectic; the $\mathrm{GL}(n)$
Poisson $\theta(Y^{-1})=\det(Y)^{1/2}\theta(Y)$ to $10^{-16}$ ($n=2,3,4$), with the elementary clocks
$E_{ij}$ landing $200/200$ random words in $\mathrm{SL}(n,\mathbb Z)$. A discrimination control
(\texttt{delta\_discriminate.py}) confirms the exact-modularity test is genuine: the true $\Delta$
theta passes at $4.5\times10^{-10}$ while reversed, sign-scrambled, and magnitude-scrambled
coefficients all fail at relative residual $1$.

\item \textbf{Local self-duality and unitary balance (Lemma~\ref{lem:selfdual}), \texttt{sym\_selfdual\_check.py}.}
For $r=0,\dots,13$: $\det\Sym^r\operatorname{diag}(\alpha,\beta)=1$ when $\alpha\beta=1$ (to $5\times10^{-15}$),
the eigenvalue multiset $\{\alpha^{r-2k}\}$ is closed under inversion (exact), and $\sum_{k=0}^r(r-2k)=0$;
the tensor identity $\det(A\otimes B)=(\det A)^{\dim B}(\det B)^{\dim A}$ holds to $3\times10^{-15}$ for
determinant-one blocks of sizes $(6,4),(7,5),(3,2)$.

\item \textbf{Symmetric powers and root numbers (\S\ref{sec:sym5}), \texttt{sym\_r\_fe2.py},
\texttt{sym\_r\_fe\_2clock.py}, \texttt{sym\_r\_sweep2.py}.} The theta self-duality
\eqref{eq:thetaselfdual} for $L(\Sym^r\Delta)$ closes with $\varepsilon=+1,-1,-1$ at $r=1,3,5$; the
two-clock Bessel completion makes the closed-form cases ($\Sym^1,\Sym^3$) exact to $3\times10^{-30}$
against a grid floor $10^{-7}$; the root numbers match the closed form
$\varepsilon(\Sym^r\Delta)=i^{-\kappa(11\kappa+1)}$ ($\kappa=\lceil(r{+}1)/2\rceil$), giving
$+,-,-,+,+,-,-$ for $r=1,\dots,13$, agreeing with the numerics for $r\le7$.

\item \textbf{Representation-agnosticism (\S\ref{sec:general}), \texttt{rep\_agnostic.py},
\texttt{three\_weight\_clock.py}, \texttt{sym\_maass.py}.} The same completion closes the functional
equation for the \emph{non-Chebyshev} Rankin--Selberg weight systems $L(\Delta\times f_{16})$,
$L(\Delta\times f_{18})$, $L(f_{16}\times f_{18})$ (weights $\{\pm13,\pm2\},\{\pm14,\pm3\},\{\pm16,\pm1\}$)
to $\sim10^{-30}$, $\varepsilon=+1$. The three-weight clock equals the two-clock composed with a third,
$(\text{Bessel}_{\mu_1,\mu_2})\star_{\mathrm M}\text{clock}_{\mu_3}$, agreeing with the direct $3$-fold
convolution to $1.7\times10^{-21}$ at arbitrary shifts. For the Maass case the two-clock goes to
imaginary order: $(2x^{it}e^{-\pi x^2})\star_{\mathrm M}(2x^{-it}e^{-\pi x^2})=4K_{it}(2\pi x)$ (the
Whittaker function) to $5\times10^{-29}$, with Mellin $\Gamma_{\mathbb R}(s+it)\Gamma_{\mathbb R}(s-it)$
to $8\times10^{-27}$, at the spectral parameter $t=9.5337$ of the first $\mathrm{SL}(2,\mathbb Z)$ Maass
cusp form.

\item \textbf{Ramified local data: the filtration tower (\S\ref{sec:general}), \texttt{ramified\_local.py},
\texttt{wild\_eps\_clock.py}.} The tame local root number is the inertia-clock Gauss sum: for primitive
$\chi$ mod $q$, $|\tau(\chi)|^2=q$ and $|\varepsilon|=1$ exactly. The wild $\varepsilon$-factor is the
higher-filtration Gauss sum over $(\mathbb Z/p^a)^*$: $|\tau|^2=p^a$ and $|\varepsilon|=1$ to
$1.7\times10^{-15}$ for $(p,a)\in\{(3,2),(3,3),(5,2),(7,2),(5,3),(11,2)\}$ (Swan $=a-1$); and the
stationary-phase tower reduction $\tau_{\mathrm{direct}}=\tau_{\mathrm{tower}}$ (the level-$2$ character
$c$ from $\chi(1+p)$, stationary $a_0\equiv-c$) holds to $3.8\times10^{-15}$ across $p=3,\dots,13$. The
monodromy $N$-ladder reduction $L_v=\det(1-\mathrm{Frob}\,q^{-s}\mid\ker N)^{-1}$ is the degree drop of
the Steinberg string.

\item \textbf{The parallel-dimension clock (\S\ref{sec:general}), \texttt{parallel\_clock.py},
\texttt{parallel\_dogfood.py}.} Quadratic base change $\Delta\mapsto\Delta_E$ over $E=\mathbb Q(i)$: the
reflected degree-four $L(\Delta_E)=L(\Delta)L(\Delta\otimes\chi_{-4})$ reproduces the base-change Satake
exactly (split/inert, diff $0$) and closes its functional equation on the equal-shift two-clock $K_0$ to
$9\times10^{-27}$, $\varepsilon=-1$. Composing clocks: $\mathrm{base\text{-}change}\circ\Sym^3$ gives a
degree-eight $L((\Sym^3\Delta)_E)$ whose two factors close on the Sym-cube Bessel $K_{11}$ ($\varepsilon=-1$
each, to $2\times10^{-31}$ and $3\times10^{-6}$).

\item \textbf{Signature-based cycle detection (\S\ref{sec:general}), \texttt{hodge\_detect.py}.} CM
algebraicity: $y^2=x^3-x$ (CM by $\mathbb Q(i)$) has $a_p=0$ at $399/399$ inert primes ($p\equiv3$ mod $4$),
against $4/399$ for the non-CM $37a$. Twist-correspondence: for $g=f\otimes\chi_5$ presented as raw
coefficients, the DC-residue slope of $L(f\times g\otimes\chi_d)$ over $d\in\{1,-4,5,8,-3,\dots\}$ reads
$0.71$ at $d=5$ and $\le0.04$ at every other $d$, with an unrelated control ($37a$) $\le0.08$ throughout ---
the hidden Tate class surfaced dynamically, no false positive. Recursive tower: the DC-residue slope of
$L(\Sym^r f\times\Sym^r f)$ for $r=1,\dots,4$ is flat ($0.90,0.90,0.82,0.81$) for the non-CM $37a$ and
climbs ($0.75,1.50,1.35,2.01$) for the CM $y^2=x^3-x$, identically ($0.77,1.51,1.36,2.03$) for
$y^2=x^3+1$ over $\mathbb Q(\sqrt{-3})$; the CM/non-CM ratio tracks the predicted pole-order tower
$1,2,2,3=\lceil(r{+}1)/2\rceil$.

\item \textbf{Loss-ledger\,$\leftrightarrow$\,cycle-filtration calibration (\S\ref{sec:general},
\S\ref{sec:status}), \texttt{hodge\_layer\_calib.py}} (from \texttt{sha\_hinge.py}, \texttt{jet\_census.py},
oracle-free). The leading central jet $L^{(r)}(E,1)/r!=\Omega\,\mathrm{Reg}\,|\Sha|\prod_p c_p/T^2$ resolves
three invariants independently: rank-$0$ curves land $|\Sha|=1$ ($11a,14a$), $4$ ($571a,960d$), $9$
($681b,2849a$) as exact squares at fixed order; the tower $37a,389a,5077a$ lands order $1,2,3$ with
$\mathrm{Reg}=0.0511,0.1524,0.4171$ and $|\Sha|=1$ (machine precision, known-value cross-check $\le
2\times10^{-5}$). Order (vanishing multiplicity), $\mathrm{Reg}$ (Abel--Jacobi datum), and $|\Sha|$
(obstruction) vary independently --- three distinct retained coordinates. The general Griffiths group is
not exhausted here; the present probe tests the induction route on homologically-trivial and delayed-depth
controls: \texttt{hodge\_griffiths\_probe.py} records the route
--- the value channel fires for every homologically-trivial cycle in the census ($3/3$ positive-order
curves, $0$ regulator-silent), and the Ceresa/Griffiths class is induced at the tensor-cube level as the
triple-product central $L$-derivative (Gross--Schoen height, Zhang \cite{ZhangGS}/Yuan--Zhang--Zhang
\cite{YZZ}) --- with the from-scratch landing deferred to the companion paper.

\item \textbf{Delayed tower signature (\S\ref{sec:status}), \texttt{tower\_exhaustion\_test.py}}
(from point-counting, oracle-free). The excess pole order of $L(\Sym^r f\times\Sym^r f)$ of the CM
$y^2=x^3-x$ over the non-CM $37a$ baseline is $-0.18$ at depth $r=1$ (silent) and $+0.67,+0.63,+1.47$ at
$r=2,3,4$ (fires), identically for $y^2=x^3+1$ --- the $0,\neq0,\dots$ signature with first nonzero slot
at depth $2$, matching the tensor level $\Sym^2$ that hosts the CM self-correspondence. Every concrete
class tested fires at a finite first depth (Tate/Mordell--Weil/$|\Sha|$ at $1$, CM at $2$, Ceresa at $3$);
none silent-forever.
\end{itemize}

\begin{thebibliography}{9}
\bibitem{AltI} S.~A.~Altu\u{g}, \emph{Beyond endoscopy via the trace formula, I}, Compositio Math.
\textbf{151} (2015).
\bibitem{AltII} S.~A.~Altu\u{g}, \emph{Beyond endoscopy via the trace formula, II}, Amer.\ J.\ Math.
\textbf{139} (2017).
\bibitem{AltIII} S.~A.~Altu\u{g}, \emph{Beyond endoscopy via the trace formula, III}, J.\ Inst.\
Math.\ Jussieu \textbf{19} (2020).
\bibitem{GJ} S.~Gelbart and H.~Jacquet, \emph{A relation between automorphic representations of
$\mathrm{GL}(2)$ and $\mathrm{GL}(3)$}, Ann.\ Sci.\ \'ENS \textbf{11} (1978).
\bibitem{Sar} P.~Sarnak, \emph{Comments on Robert Langlands' lecture ``Endoscopy and Beyond''}, notes
(2001), \texttt{publications.ias.edu/sarnak}.
\bibitem{Hormander} L.~H\"ormander, \emph{The Analysis of Linear Partial Differential Operators I},
2nd ed., Grundlehren \textbf{256}, Springer (1990); Thm.~7.7.1 (non-stationary phase).
\bibitem{Wong} R.~Wong, \emph{Asymptotic Approximations of Integrals}, Classics in Applied
Mathematics \textbf{34}, SIAM (2001); Ch.~3 (Watson's lemma and the Mellin method).
\bibitem{Tate} J.~Tate, \emph{Fourier analysis in number fields and Hecke's zeta-functions}, thesis
(1950), in J.~W.~S.~Cassels and A.~Fr\"ohlich (eds.), \emph{Algebraic Number Theory}, Academic Press
(1967), 305--347.
\bibitem{GJa} R.~Godement and H.~Jacquet, \emph{Zeta Functions of Simple Algebras}, Lecture Notes in
Mathematics \textbf{260}, Springer (1972).
\bibitem{CPS} J.~W.~Cogdell and I.~I.~Piatetski-Shapiro, \emph{Converse theorems for
$\mathrm{GL}_n$}, Publ.\ Math.\ IH\'ES \textbf{79} (1994), 157--214; \emph{II}, J.\ reine angew.\
Math.\ \textbf{507} (1999), 165--188.
\bibitem{Shahidi} F.~Shahidi, \emph{Eisenstein Series and Automorphic $L$-functions}, Amer.\ Math.\
Soc.\ Colloquium Publications \textbf{58}, AMS (2010); and \emph{A proof of Langlands' conjecture on
Plancherel measures; complementary series for $p$-adic groups}, Ann.\ of Math.\ \textbf{132} (1990),
273--330.
\bibitem{KimShahidi} H.~Kim and F.~Shahidi, \emph{Functorial products for
$\mathrm{GL}(2)\times\mathrm{GL}(3)$ and the symmetric cube for $\mathrm{GL}(2)$}, Ann.\ of Math.\
\textbf{155} (2002), 837--893; H.~Kim, \emph{Functoriality for the exterior square of
$\mathrm{GL}_4$ and the symmetric fourth of $\mathrm{GL}_2$}, J.\ Amer.\ Math.\ Soc.\ \textbf{16}
(2003), 139--183.
\bibitem{NT} J.~Newton and J.~A.~Thorne, \emph{Symmetric power functoriality for holomorphic modular
forms, I \& II}, Publ.\ Math.\ IH\'ES \textbf{134} (2021), 1--116, 117--152.
\bibitem{Serre} J.-P.~Serre, \emph{Abelian $\ell$-adic Representations and Elliptic Curves}, Benjamin
(1968); Ch.~I, the equidistribution criterion for the Sato--Tate law from the analyticity and
non-vanishing of the symmetric-power $L$-functions.
\bibitem{JS} H.~Jacquet and J.~A.~Shalika, \emph{A non-vanishing theorem for zeta functions of
$\mathrm{GL}_n$}, Invent.\ Math.\ \textbf{38} (1976), 1--16; non-vanishing on $\Re s=1$ from the
positivity of the Rankin--Selberg convolution; and \emph{On Euler products and the classification of
automorphic representations}, Amer.\ J.\ Math.\ \textbf{103} (1981), the local bound $|\alpha|<p^{1/2}$.
\bibitem{LRS} W.~Luo, Z.~Rudnick, and P.~Sarnak, \emph{On the generalized Ramanujan conjecture for
$\mathrm{GL}(n)$}, in \emph{Automorphic Forms, Automorphic Representations, and Arithmetic}, Proc.\
Sympos.\ Pure Math.\ \textbf{66}, Part 2, AMS (1999), 301--310; the bound
$|\alpha|\le p^{1/2-1/(n^2+1)}$ toward Ramanujan.
\bibitem{Titchmarsh} E.~C.~Titchmarsh, \emph{The Theory of Functions}, 2nd ed., Oxford University
Press (1939), \S\S5.6--5.65 (Phragm\'en--Lindel\"of).
\bibitem{ZhangGS} S.~Zhang, \emph{Gross--Schoen cycles and dualizing sheaves}, Invent.\ Math.\
\textbf{179} (2010), 1--73; the Beilinson--Bloch height of the modified diagonal.
\bibitem{YZZ} X.~Yuan, S.~Zhang, and W.~Zhang, \emph{Triple product $L$-series and Gross--Schoen
cycles}, book manuscript; the central derivative of the triple-product $L$-function as a
Gross--Schoen/Ceresa height.
\end{thebibliography}

\end{document}
